% SOURCE docs/release/2026-10-29/source/docs/writing-repair/2026-10-29/src/f50f32c01b8a03ec.tex CHAR 0
\chapter{Evaluating causal assistance without changing the question}
\label{la:evaluation}

The question throughout this chapter is whether assistance preserves a banking decision while improving its resolution at an acceptable cost. The opening defines the comparison; the coverage screen and prospective experiment then show what a finite evaluation could establish. The source studies test the reference answers, exposure history, adjudication and scoring rules on which that experiment would rely. Finally, the evaluator, memory and toolkit comparisons follow those requirements into a working interface. Their different failure mechanisms cannot be reduced to a single accuracy score.

\section{The benchmark unit is a causal problem}

An agent can produce a plausible report, execute valid code and still answer
the wrong causal question. Evaluation must follow the entire chain from
business meaning through assumptions and identification to uncertainty,
availability and decision loss. A benchmark of text fluency, graph edit
distance or successful tool calls measures only part of that chain.

CLadder demonstrates the value of formal questions with varied verbalizations,
including commonsense and deliberately unfamiliar or conflicting stories.
The inspected paper also identifies language-to-formal-tool formulation as
a continuing challenge. Its historical binary-answer accuracy is evidence
about those tasks and models; it is not a measurement of effect coverage,
decision loss or banking readiness. \cite{v12_la_2312_04350}

CausalDS expands the evaluation object to narrated synthetic mechanisms,
data and a separate observation model. Its inspected results show why
symbolic success, uncertainty quality, abstention and cost should be reported
separately. Its author-reported exam and repeated runs are useful design
evidence, not model evaluations conducted in this project. A generated world
with known truth also remains different from authentic bank data whose
semantic and assignment mechanisms must be established externally.
\cite{v12_la_2607_08093}

The proposed experimental unit is therefore an independently generated or
independently adjudicated causal mechanism and question. Several phrasings,
label conditions or stochastic model calls can be nested within that unit.
They do not create new independent worlds. Repeating one problem fifty times
can measure response variation for that problem; it does not provide fifty
independent demonstrations of general causal competence.

\section{Matched information and meaningful alternatives}

The prospective design compares a fixed numerical workflow, a language-only
workflow, retrieval assistance, a hybrid of numerical tools and typed model
proposals, and structured expert workflows with and without assistance.
Each arm needs a declared information budget. If one receives the true
causal graph and another must discover it, their effect-estimation scores
do not isolate the value of language assistance. If retrieval adds a valid
institutional fact, that information gain should be distinguished from a
better inference algorithm.

Eight strata are proposed: independent sufficient DAGs; latent-variable
graph classes; temporal data with known clocks; treatment-confounder
feedback; panel or quasi-experimental assignment; proxies, linkage and
selection; weak or nonidentifying effects; and ambiguous business targets
or monetary loss. Label and evidence conditions cross ordinary descriptions,
unfamiliar labels, misleading or irrelevant material, and conflict between
numerical evidence and text. These strata test different failure mechanisms
rather than treating all causal questions as interchangeable.

The current planning allocation is 400 mechanisms, with 50 per stratum and
three stochastic repetitions nested within each mechanism. This is an
allocation for developing the design, not a frozen final sample size or a
power calculation. For orientation, the worst-case normal approximation for
one Bernoulli proportion at 400 independent observations has half-width
\(1.96\sqrt{0.25/400}\), approximately 0.049. Paired differences, stratified
outcomes and rare severe errors require their own precision design. A
development pilot would inform that design and must be excluded from the
final test. The six task families developed in Section~\ref{exec13:live-design} describe workflow functions, whereas the eight strata here describe scientific conditions; they are different design dimensions. That later section uses additional sample sizes to illustrate error control, power and cost. Those calculations refine the choice of a trial rather than combine into an approved allocation.

Both natural-workflow and matched-budget comparisons can be useful. The
first asks what an analyst can achieve with a realistic workflow, including
review. The second asks what changes when information or computation is
controlled. Neither should hide extra expert time, repeated retrieval or
retry costs. The comparison must state which question it answers.

\section{A target error can survive perfect arithmetic}

An exact project IV construction illustrates a benchmark-label hazard.
Let a randomized binary instrument \(Z\) affect treatment for half the
population, the compliers; the other half never receives treatment. Set
\(Y(0)=0\) for everyone, \(Y(1)=1\) for compliers and \(Y(1)=3\) for
never-takers. Treatment equals \(Z\) for compliers and zero otherwise.
Randomization, exclusion and monotonicity hold in this construction.

The observed outcome difference across instrument values is \(0.5\) and
the treatment-probability difference is \(0.5\), so the Wald ratio is one.
The population ATE is two. Changing the never-takers' unobserved treated
outcome changes the population ATE without changing the observed law. A
scorer that calls the unrestricted Wald ratio the population ATE would
reward a target substitution. This is a finite counterexample supporting
review of generator assumptions, not a claim that every IV question in a
named benchmark uses this construction.

The same concern applies to trimming, overlap weighting and selected
populations. A language answer can reproduce the right formula for a
different target. Evaluation should therefore compare the population,
measure, action, comparator, horizon and units before comparing the numeric
answer. Critical semantic errors deserve their own count, even when a
numeric discrepancy happens to be small in the test mechanism.

\section{Graphs, sets and refusals need separate scores}

Graph evaluation should compare adjacency and endpoint statements against
the appropriate identifiable class. A justified circle or undirected edge
should not be penalized merely because a simulator internally selected one
fully directed graph. Orientation informativeness can be reported alongside
validity. This makes a method that preserves genuine ambiguity distinguishable
from one that returns little information and from one that confidently
overorients the graph.

Identification outcomes also need separate labels. Scientific
nonidentification, unsupported implementation, malformed output, explicit
abstention, timeout and missing input are not the same event. A well-designed
agent can correctly identify a missing premise without successfully computing
an effect. A backend can be unable to implement a valid functional. Combining
these into one ``incorrect'' or ``no answer'' category would obscure the
remedy.

For identified scalar targets, bias and error can be meaningful. For reported
sets, evaluate coverage, topology, finite availability and size. Whole-line
answers can have high coverage and little decision value. A method that
returns short intervals on a selected easy subset needs an all-request
availability denominator. Unbounded and disconnected sets must remain in the
evaluation rather than being replaced by finite plotting limits.

For an eligible finite interval \([l,u]\), one prospective score at nominal
level \(1-\alpha\) is
\[
 (u-l)+\frac{2}{\alpha}(l-\theta)1\{\theta<l\}
       +\frac{2}{\alpha}(\theta-u)1\{\theta>u\}.
\]
The formula balances length with penalties for missing the target. It is
not an instruction to compress a union of intervals into one interval. For
unions, coverage and total Lebesgue length can be reported while retaining
infinite lengths. Raw losses should not be capped simply to make a plot or
ranking convenient. The eligible target type must be stated before the
score is selected.

\section{Contamination, adaptation and selective reporting}

An exposed demonstration cannot become a held-out evaluation merely because
its variables are renamed. The local problem fixtures and replay examples
have already been used to develop and debug the prototype. They can continue
to test conformance, but they cannot provide independent estimates of a
language model's causal accuracy.

The final design needs a frozen development/test split, mechanism identities,
prompts, model/provider revisions, retrieval corpus and ranking rules, numerical
baselines, stopping rules and analysis code. Mechanism families should be
separated where family reuse would reveal the answer. Label novelty and
mechanism novelty should be varied independently. An unfamiliar name attached
to a memorized mechanism is not the same test as an unfamiliar mechanism.

Test outcomes cannot be used to tune prompt templates, choose the best
retriever, strengthen selected priors or discard difficult cells while
retaining a final-test interpretation. A changed method requires a new
reviewed evaluation revision. Failed and unavailable requests remain in the
record. Reporting only the critic-approved subset would confound answer
quality with selection by the critic.

The inspected OCI paper provides a useful refusal example: its critic
declines many replications when overlap is poor. That can be preferable
to an unsupported estimate, but it makes the denominator essential.
The reported successful answers, refused requests, review effort and
remaining target population must all be visible. The external results
are author reports, not a validation of this project's critic.
\cite{v12_la_2607_22443}

\section{Decision value and cost}

The ultimate question is whether assistance improves feasible decisions at
an acceptable total cost. In a known synthetic mechanism, a policy's expected
loss can be calculated under the mechanism and compared with the chosen
baseline. In a bank, the feasible policy, valuation model, risk limits,
operating costs and outcome truth need authentic inputs. Synthetic decision
loss does not establish institutional return on investment.

The prospective cost record includes model calls, input/output/cache tokens,
retrieval, numerical CPU time, wall time, memory, retries and human review
and rework. Actual dated prices are needed to turn these quantities into a
currency total. Unknown prices remain missing; zero live calls in an offline
prototype do not imply a free live workflow. A fast answer that requires
substantial correction can have a worse total cost than a slower transparent
analysis.

No final live-model experiment is reported by this chapter. Five prerequisite
classes remain: an exact live backend and custody terms; actual prices and
a bounded budget; independent held-out mechanisms and original benchmark
bytes; named business/diagram adjudicators and a rubric; and qualified
numerical baselines. These are missing program inputs. No credentials were
read and no provider endpoint was probed to infer that a service is unavailable.
The design provides a route to evaluation once those inputs are supplied;
it provides no current empirical model-performance claim. The remaining studies turn these requirements into inspectable tests: first the probability of passing a coverage rule, then a prospective comparison, and finally the reliability of its references, evaluators and numerical handoffs.


% SOURCE docs/release/2026-10-29/source/docs/writing-repair/2026-10-29/src/577491ab16aa8ade.tex CHAR 0
\section{The probability of passing a coverage screen}
\label{v13:la-screen}

Evaluating an agent that recommends inferential procedures requires a clear
distinction between nominal coverage, an argument for coverage and an empirical
qualification rule. Suppose a fixed procedure is tested in $n$ independently
generated replications of a declared mechanism. Let $K$ count requests whose
reported sets cover the truth; unavailable or failed requests remain in the
declared denominator. If the per-request coverage is $p$, then under these
simulation assumptions $K$ has a binomial law.

For a simultaneous family of $m$ cells and total error budget $\alpha$, consider
the one-sided exact lower bound
\[
 L(K,n)=\operatorname{BetaQuantile}
       (\alpha/m;K,n-K+1),\qquad K>0,
\]
with $L(0,n)=0$. A rule requiring $L\geq0.95$ certifies a minimum in this
specific finite experiment. At true $p=0.95$, its probability of passing one
cell is at most $\alpha/m$. This is how a lower-bound certification test behaves
at its boundary, not evidence that the procedure's true coverage is zero.

Before qualification, a designer can compute the smallest integer $k_*$ that
passes and the operating characteristic
\[
 \Pr_p(\hbox{pass})=\Pr_p(K\geq k_*).
\]
The project calculation evaluates this exactly for several replication counts,
family sizes and values of $p$, and independently checks the equivalent
binomial-tail threshold. It draws no random data and does not test a causal
estimator. Its purpose is to make the qualification design interpretable.
\cite{v13_project_exact}

Increasing replications can distinguish a genuine coverage margin above 95\%
more sharply. It cannot make a level-controlled test have high power at the
exact boundary. This does not justify lowering the required minimum after
seeing results. The analyst must state what guarantee is theoretically justified,
what empirical screen is being applied and what its power is. A method may
have a plausible asymptotic argument while failing this finite screen; that
should be reported as a qualification shortfall rather than silently recoded
as success.

An agent could otherwise exploit this distinction in two bad ways. It might
declare a method invalid merely because a low-powered boundary screen failed,
or widen intervals until they pass while ignoring decision usefulness. The
evaluation should penalize both. Coverage, finite availability and useful width
are separate dimensions. A whole-line interval covers every finite truth but
may support no action. Bank width and loss thresholds require actual decision
owners; synthetic thresholds should remain labeled as such.

\section{A live comparison needs independent mechanisms and complete costs}

An offline controller can demonstrate request validation, abstention, tool
dispatch and preservation of graph/set semantics. It cannot measure how a live
model translates a new banking problem, how often retrieval changes its graph,
or what the workflow costs. Zero model calls means those empirical quantities
are unmeasured; it does not mean accuracy is zero or costs are free.

A credible live comparison binds the exact backend/revision and inference
settings, data-custody arrangement, actual prices and approved budget, independent
held-out mechanisms or benchmark bytes, adjudicators and rubric, and qualified
numerical baselines. The final prompts, retrieval corpus, tools, retry rules,
token/time limits and numerical tolerances should be frozen before the first
final request. Development cases already exposed to the agent cannot become
held-out by receiving different filenames.

Compare data-only, model-only, retrieval-assisted, hybrid and expert-assisted
workflows under matched information and budgets. A hybrid receives additional
information when it uses a valid domain exclusion; its improvement should not
be attributed solely to reasoning. Conversely, a data-only baseline should not
be asked to identify an effect that the observed law cannot identify. The
equal-law counterexamples in Section~\ref{v13:la-equivalence} make this boundary
testable.

Metrics should include whether the translated intervention/population is valid,
whether graph output respects equivalence and latent ambiguity, whether the
downstream effect is scientifically supported, and whether abstention occurs
for the right reason. Repeated-run disagreement is relevant even when the
average score is favorable. Fluent explanations and exact string matches to
one graph are inadequate substitutes for these criteria.

Complete costs include unsuccessful calls, retries, retrieval, tool computation,
human review and operating overhead. If an agent frequently produces a request
that needs manual correction, its low token bill can conceal high workflow cost.
Decision value should be evaluated on feasible actions and held-out outcomes
or an explicitly justified reference, with uncertainty. A stipulated arithmetic
example explains the accounting but cannot establish provider pricing, model
calibration or bank ROI.

\section{Testing semantic failures before the expensive study}

Several low-cost deterministic cases can be exercised before live qualification.
A y0 nonidentification response must not become a zero effect. A continuous
density must not be interpreted as a binary probability. A grouped allocation
contrast must not become an individual direct effect. A statsmodels internal
covariance option must not be reported as final clustered uncertainty when the
inspected implementation supplies HC0. \cite{v13_sw_y0,v13_sw_causalgps,v13_sw_inferference,v13_sw_statsmodels}

These cases test the proposed interface and evaluation rubric. They are useful
development evidence, and their exposure should be recorded. Fresh independent
mechanisms are still necessary to evaluate generalization. Keeping those roles
separate allows an agentic workflow to improve through development while
retaining a meaningful final assessment.


% SOURCE docs/release/2026-10-29/source/docs/writing-repair/2026-10-29/src/dcc5dec0e2bcb08d.tex CHAR 0
\hypertarget{evaluate-the-claim-behind-the-benchmark-score}{%
\subsection{Evaluate the claim behind the benchmark
score}\label{evaluate-the-claim-behind-the-benchmark-score}}

\label{exec13:census-evaluation}

CausalEval places diverse language tasks in one comparative survey,
including commonsense cause selection, correlation-to-causation
questions, graph-based inference and judgments about described events.
This is useful navigation, but the common label ``causal reasoning''
does not make their outcomes interchangeable. Correctly classifying a
relation mentioned in a sentence, recovering a graph from observations
and identifying a policy value test different information and
assumptions. \cite{exec13_census_causaleval}

\begin{longtable}[]{@{}lll@{}}
\toprule
\begin{minipage}[b]{0.30\columnwidth}\raggedright
Evidence object\strut
\end{minipage} & \begin{minipage}[b]{0.30\columnwidth}\raggedright
What the inspected source measures\strut
\end{minipage} & \begin{minipage}[b]{0.30\columnwidth}\raggedright
Additional evidence needed for the proposed causal role\strut
\end{minipage}\tabularnewline
\midrule
\endhead
\begin{minipage}[t]{0.30\columnwidth}\raggedright
Susanti--Holsmoelle relation extraction\strut
\end{minipage} & \begin{minipage}[t]{0.30\columnwidth}\raggedright
Binary relations between entities in text\strut
\end{minipage} & \begin{minipage}[t]{0.30\columnwidth}\raggedright
Joint graph constraints, target-population meaning and effect
identification\strut
\end{minipage}\tabularnewline
\begin{minipage}[t]{0.30\columnwidth}\raggedright
Sng and colleagues' hallucination study\strut
\end{minipage} & \begin{minipage}[t]{0.30\columnwidth}\raggedright
Consistency of causal ratings, including a supplied mocked graph\strut
\end{minipage} & \begin{minipage}[t]{0.30\columnwidth}\raggedright
Accuracy against independently justified causal truth\strut
\end{minipage}\tabularnewline
\begin{minipage}[t]{0.30\columnwidth}\raggedright
Augur\strut
\end{minipage} & \begin{minipage}[t]{0.30\columnwidth}\raggedright
Trend prediction and explanation consistency\strut
\end{minipage} & \begin{minipage}[t]{0.30\columnwidth}\raggedright
Intervention recovery under declared temporal and confounding
assumptions\strut
\end{minipage}\tabularnewline
\begin{minipage}[t]{0.30\columnwidth}\raggedright
ORCA and CausalPulse\strut
\end{minipage} & \begin{minipage}[t]{0.30\columnwidth}\raggedright
Workflow outcomes plus distinct scientific metrics\strut
\end{minipage} & \begin{minipage}[t]{0.30\columnwidth}\raggedright
Defined oracles, retained failures and matched target-level
evaluation\strut
\end{minipage}\tabularnewline
\begin{minipage}[t]{0.30\columnwidth}\raggedright
CausalSmith\strut
\end{minipage} & \begin{minipage}[t]{0.30\columnwidth}\raggedright
Reported proof checks and internal run dispositions\strut
\end{minipage} & \begin{minipage}[t]{0.30\columnwidth}\raggedright
Statement fidelity, independent scientific significance and target
assumptions\strut
\end{minipage}\tabularnewline
\bottomrule
\end{longtable}

\hypertarget{relation-extraction-is-a-serious-baseline-with-a-narrower-claim}{%
\subsubsection{Relation extraction is a serious baseline, with a
narrower
claim}\label{relation-extraction-is-a-serious-baseline-with-a-narrower-claim}}

Susanti and Holsmoelle compare prompting with fine-tuned models on four
relation datasets. On DDI, Table 1 reports F1 of 67.4 for Single-Prompt
C and 87.9 for PubMedBERT: a 20.5-point difference. The paper reports
five-fold averages and explicitly limits its evaluation to binary
pairwise causality assessment. It neither demonstrates joint graph
verification nor establishes a universal advantage for fine-tuning under
equal annotation costs. The comparison does establish that a supervised
smaller model deserves consideration when the proposed task is
extraction from a stable document domain.
\cite{exec13_census_graph_validation}

For banking notes, evaluate the extraction role on held-out documents
with independently specified labels and timestamps. Track abstention and
the treatment of ambiguous or contradictory passages. Then separately
evaluate whether extracted statements are admissible graph evidence. A
sentence that says ``the waiver caused repayment'' is evidence of what
the author claimed; it is not, by itself, evidence that assignment was
unconfounded. A useful agent should retain the quotation, its source and
the alternative causal interpretations.

Sng and colleagues illustrate why agreement is a separate target. Their
retrieval experiment supplies a mocked true causal graph, and their
conclusion explicitly says that reducing inconsistent ratings does not
imply accurate causal prediction. A comparison in which one arm receives
such a graph evaluates its use of supplied information. It does not
measure whether retrieval can discover the graph in a bank setting.
Registering both correct and misleading retrieved evidence is necessary
to expose that distinction. \cite{exec13_hallucination}

\hypertarget{task-specific-gains-should-remain-task-specific}{%
\subsubsection{Task-specific gains should remain
task-specific}\label{task-specific-gains-should-remain-task-specific}}

In CausalEval's Table 3, GPT-4o scores 69 on CLADDER with ordinary
chain-of-thought and 76 with FunctionCall, but its e-CARE scores are 90
and 73, respectively. These are author-reported pass@1 results for the
named versions and samples. The gains and losses should remain visible
rather than being replaced by the claim that tools always improve causal
reasoning. The table is also preferable to silently repeating
inconsistent numerical values in the paper's surrounding discussion.
\cite{exec13_census_causaleval}

The study repeats each question five times with greedy decoding.
Repeated calls can reveal response variation, but five calls to the same
question are not five independent causal problems. For illustration,
suppose there are \(N\) independent mechanisms, \(m\) calls per
mechanism, equal marginal score variance \(\sigma^2\), and common
within-mechanism correlation \(\rho\). Then

\[
 \operatorname{Var}(\bar S)
 =\frac{\sigma^2}{Nm}\{1+(m-1)\rho\}.
\]

At \(N=100,m=5,\rho=1\), the effective independent count is 100, not
500. This is a project-derived sampling illustration, not an estimate of
the paper's unreported correlation. The appropriate prototype analysis
clusters calls and paraphrases by underlying mechanism, reports the
number of distinct mechanisms, and uses paired comparisons when methods
see the same cases.

CausalEval also evaluates some expensive models on 100 cases rather than
the usual 500. Its human row combines scores from original benchmark
papers with a study of three graduate students on 100 questions for
remaining datasets. A combined human/model gap therefore is not one
uniform matched experiment. Any investment argument should state the
comparison actually performed, the cost and failure accounting, and the
uncertainty that remains.

\hypertarget{an-affordable-falsifiable-evaluation-sequence}{%
\subsubsection{An affordable, falsifiable evaluation
sequence}\label{an-affordable-falsifiable-evaluation-sequence}}

The proposed sequence starts with fixed causal mechanisms, explicit
public document provenance and deterministic numerical answers. It then
tests formulation, variable timing, graph uncertainty, identification,
execution and reporting separately, followed by their composition.
Baselines receive the same data, permitted documentation and tools.
Model-only reasoning, retrieval, a typed numerical workflow, and a
supervised extractor answer different engineering questions and should
be included where relevant.

The primary denominator includes every attempted case, with scientific
refusals distinguished from crashes and wrong answers. Numerical
accuracy conditional on successful execution is a secondary outcome.
Cost includes failed calls, retrieval, annotation, retries and human
review. Deliberately adverse cases include an accurate quotation
supporting a false causal claim, a timestamp that makes a proposed
control post-treatment, alternative graphs with different effects, a
misleading retrieved graph, and an unsupported edit policy.

Public and generated data can establish that these failure cases are
handled as specified and can compare methods under a known law. They do
not establish the frequency of those failures in a target bank or the
value of changing its real decisions. The next expensive experiment
should address a remaining uncertainty that could change the build or
adoption decision. Until such evidence exists, the proposal can
recommend a bounded prototype and a testable research direction while
preserving uncertainty about live-model performance, source completeness
and bank transport.

\label{exec13:census-evaluation-end}


% SOURCE docs/release/2026-10-29/source/docs/writing-repair/2026-10-29/src/a061406f90d833b5.tex CHAR 0
\hypertarget{a-prospective-experiment-for-causal-assistance}{%
\subsection{A prospective experiment for causal
assistance}\label{a-prospective-experiment-for-causal-assistance}}

\label{exec13:live-design}

The first live comparison should determine whether language assistance
improves a specified causal workflow enough to justify its complete
cost. A fluent explanation, a correctly guessed graph and a correctly
identified effect are different outcomes. The design therefore evaluates
formulation, information timing, graph ambiguity, identification,
numerical execution and reporting as separate tasks, followed by their
composition. Public documents and generated mechanisms can support this
experiment before a bank funds access to proprietary records. Model
revisions, final sample sizes, acceptance thresholds and prices remain
to be selected from actual inputs; the calculations below are planning
examples, not a registered trial or evidence of live-model performance.

CausalEval's inspected experiment repeats each question five times and
averages the results within question. Such repeats can expose output
variation, but cannot create additional independent causal problems
\citep{exec13_census_causaleval}. The same issue arises in the
proposal's authored tasks. Its development collection contains 24
variants in twelve template groups; its reserved collection contains 48
variants in twenty-four groups, with two variants in every group. These
assets support interface and rubric development. Their authorship and
generation history do not establish independent, identically distributed
final mechanisms or independently adjudicated causal truth
\citep{exec13_live_planning}.

\hypertarget{mechanisms-information-and-the-reference-answer}{%
\subsubsection{Mechanisms, information and the reference
answer}\label{mechanisms-information-and-the-reference-answer}}

A final task should begin with a declared scientific information state.
For a generated structural model, the evaluator can retain the full
graph, assignment mechanism and intervention outcomes while the analysis
sees only its permitted observations and metadata. The reference must
answer two questions: what is true in the generating world, and what is
identified from the information supplied to the method? They need not
have the same answer. An agent that guesses the hidden true effect in a
nonidentifying problem has not correctly solved identification.

The task portfolio retains six functional families: business
formulation; clocks and selection; graph construction and discovery;
identification and estimator choice; evidence tracing; and tool
composition. Across them, stress conditions include hidden confounding,
temporal dependence, treatment-- confounder feedback, panel assignment,
imperfect linkage and ambiguous monetary objectives. A selected first
experiment need not claim every condition at once. Its included cells,
excluded capabilities and substantive alternatives should be fixed
before evaluation rather than chosen from successful outputs.

The split should occur at mechanism and template-family level before
paraphrases, renamed variables or translations are created. A new name
for a familiar graph does not make its mechanism unseen. Public
benchmarks require their own provenance and label review; public
availability does not establish absence from a model's training data. A
generated final collection must have its generator, information
restrictions and scoring rules checked before the final draw. Prompts
and retrieval settings are chosen on a separate development collection.
Any inspected final case is recorded as exposed, and subsequent method
changes require a new evaluation revision.

Author-reviewed algebraic references can support a bounded proposal
demonstration. Ambiguous business language, source applicability and
multilingual meaning may also require human judgment. Such judgments
should use a frozen rubric, method-blinded presentation and retained
disagreements. Until independent adjudicators actually participate, the
result must remain an author-reviewed study. Neither a generated label
nor agreement between two model responses substitutes for authentic bank
assignment evidence.

The adverse conditions should change both language and information.
Useful paired variants include an unfamiliar variable vocabulary, a unit
change, a delayed feature timestamp, a quoted but scientifically
unsupported causal claim, a misleading retrieved graph, and conflict
between text and observed data. Additional languages require a checked
translation that preserves the population, intervention, clock and
uncertainty. Language versions and stochastic repeats stay nested within
the original mechanism. At least one primary variant must be chosen
prospectively for the binary comparison derived below; sensitivity
results on the other variants remain visible.

\hypertarget{what-each-comparator-is-allowed-to-know}{%
\subsubsection{What each comparator is allowed to
know}\label{what-each-comparator-is-allowed-to-know}}

The comparison should hold a backend fixed while changing its
information and tool access. This separates the value of retrieval and
numerical composition from a change in the underlying model. Comparing
backends can be a later, explicitly costed dimension. Four proposed live
arms are a frozen prompt alone; the same model with audited retrieval;
the same model with typed numerical tools; and a hybrid with both. A
critique step can be ablated within the hybrid when it fits the same
declared total resource budget. It is not free additional computation.

The model-only arm receives the same declared task and permitted
summaries, but no hidden graph or numerical-tool result. Retrieval arms
use an immutable corpus and a fixed query/ranking policy, and save the
returned passages. Tool arms share pinned algorithms and the same
observations. Constraint proposals must state their evidence and can be
refused by the numerical operation. A proposed restriction cannot
silently change the population, sampling law or model class to make an
inconvenient query identifiable.

Non-model comparators remain essential. A deterministic form/schema
workflow tests whether structured elicitation is sufficient. A data-only
estimator tests the numerical stage on admissible structured inputs. A
retrieval-only system tests whether locating a primary passage suffices.
These comparators have different roles: a data-only estimator given a
correct structured estimand is not simultaneously solving free-text
business translation. Report stage-specific results and end-to-end
results with the information differences exposed. Human and
human-plus-model workflows are optional additional arms whose time,
expertise and reference independence must be measured rather than
inferred from a model benchmark.

Equal budgets should cover total tokens, model calls, tool calls,
elapsed time and retries per assigned task. If one arm consumes less
than its cap, retain that advantage in the cost account. If an agent
uses more calls, compare its benefit against the additional actual cost.
A transport retry and a deliberate internal critique are different
events; both belong in the call ledger. Timeouts, parse failures and
exhausted budgets remain assigned outcomes.

\hypertarget{useful-abstention-and-unjustified-certainty}{%
\subsubsection{Useful abstention and unjustified
certainty}\label{useful-abstention-and-unjustified-certainty}}

Let \(U_i\) indicate a useful, scientifically correct resolution of an
assigned task, \(E_i\) a critical scientific error, and \(A_i\)
availability of a usable response. The rubric requires \(U_i\leq A_i\)
and \(U_iE_i=0\). A useful resolution can be a justified effect set or a
specific request for missing information. It cannot be an unsupported
point estimate or a generic refusal unrelated to the task.

\begin{longtable}[]{@{}lll@{}}
\toprule
\begin{minipage}[b]{0.30\columnwidth}\raggedright
Reference state\strut
\end{minipage} & \begin{minipage}[b]{0.30\columnwidth}\raggedright
Useful response\strut
\end{minipage} & \begin{minipage}[b]{0.30\columnwidth}\raggedright
Outcome that must remain a failure\strut
\end{minipage}\tabularnewline
\midrule
\endhead
\begin{minipage}[t]{0.30\columnwidth}\raggedright
Identified target with supported computation\strut
\end{minipage} & \begin{minipage}[t]{0.30\columnwidth}\raggedright
Correct estimand, eligible method and traceable result\strut
\end{minipage} & \begin{minipage}[t]{0.30\columnwidth}\raggedright
Blanket abstention or a result for another population\strut
\end{minipage}\tabularnewline
\begin{minipage}[t]{0.30\columnwidth}\raggedright
Several compatible effects\strut
\end{minipage} & \begin{minipage}[t]{0.30\columnwidth}\raggedright
Preserve the admissible ambiguity and explain which information would
narrow it\strut
\end{minipage} & \begin{minipage}[t]{0.30\columnwidth}\raggedright
Select a favorable graph and report its effect as certain\strut
\end{minipage}\tabularnewline
\begin{minipage}[t]{0.30\columnwidth}\raggedright
Missing decision clock or action definition\strut
\end{minipage} & \begin{minipage}[t]{0.30\columnwidth}\raggedright
Ask the precise question whose answer changes the analysis\strut
\end{minipage} & \begin{minipage}[t]{0.30\columnwidth}\raggedright
Invent a clock, treatment or business objective\strut
\end{minipage}\tabularnewline
\begin{minipage}[t]{0.30\columnwidth}\raggedright
Invalid or unsupported request\strut
\end{minipage} & \begin{minipage}[t]{0.30\columnwidth}\raggedright
Explain the actual limitation and a valid next route\strut
\end{minipage} & \begin{minipage}[t]{0.30\columnwidth}\raggedright
Pretend that the requested calculation was executed\strut
\end{minipage}\tabularnewline
\begin{minipage}[t]{0.30\columnwidth}\raggedright
Unavailable backend or timeout\strut
\end{minipage} & \begin{minipage}[t]{0.30\columnwidth}\raggedright
Retain the operational failure in the denominator\strut
\end{minipage} & \begin{minipage}[t]{0.30\columnwidth}\raggedright
Relabel infrastructure failure as scientific nonidentification\strut
\end{minipage}\tabularnewline
\bottomrule
\end{longtable}

Operational honesty does not automatically count as a successfully
resolved scientific task. Availability and unsupported-capability rates
are reported separately, including among identifiable tasks.
Critical-error rates among answered cases are also reported when their
denominator is nonzero. This prevents a system with no available answers
from winning because it emitted no false causal claims. An evaluation must consider useful answers, availability, critical
errors and cost together; the required numerical thresholds are not supplied by the
current proposal.

Graph scoring compares the returned object with the admitted equivalence
class. Unsupported arrowheads count separately from omitted informative
orientations, and a justified circle is not penalized as if a fully
directed DAG were identified. Effect scoring first checks the target and
its identification status. Scalar error is meaningful only for the
admitted scalar target; set results retain disconnected components,
empty sets and unbounded regions. Coverage and length are separate
outcomes, and infinite length stays infinite. Monetary policy loss is
evaluated only under the declared generated mechanism and feasible
actions. It is not a measurement of bank ROI.

\hypertarget{how-many-independent-problems-could-justify-a-claim}{%
\subsubsection{How many independent problems could justify a
claim?}\label{how-many-independent-problems-could-justify-a-claim}}

Consider an illustrative critical-error indicator for independent,
identically distributed mechanism draws in one task family. If no errors
are observed in \(n\) draws, then the probability of this event at an
error rate \(p\) is \((1-p)^n\). Inverting this probability yields the
one-sided upper confidence bound \[
 p_U=1-\alpha_*^{1/n}.
\] With a single 95\% bound, zero errors in 48 independent mechanisms
still gives \(p_U=0.06050\). If those forty-eight cases were only
twenty-four independent mechanisms with duplicated observations,
substituting 48 for 24 would overstate the evidence; the corresponding
bound at 24 is 0.11735. The authored reserved templates do not
themselves establish either sampling law \citep{exec13_live_planning}.

For \(K\) simultaneous one-sided bounds, allocating
\(\alpha_* = 0.05/K\) gives a familywise guarantee by the union bound.
To obtain \(p_U\leq0.05\) after zero errors requires \[
 n\geq\left\lceil
       \frac{\log(0.05/K)}{\log(0.95)}
       \right\rceil.
\] The resulting counts are 59 for one claim, 94 for six, and 107 for
twelve. Here 5\% is an illustrative error tolerance, not a bank
requirement. Six task-family error bounds and six improvement tests
would together need a joint allocation over twelve claims. Assigning
each collection a separate 5\% budget would not provide the same
combined guarantee. No independence between families is needed for this
multiplicity argument, although the binomial model within each family
still needs its stated sampling law.

Improvement also needs prospective power planning. On the designated
primary variant of mechanism \(i\), define
\(D_i=U_i^{\rm hybrid}-U_i^{\rm comparator}\). Let \(p_+=\Pr(D_i=1)\),
\(p_-=\Pr(D_i=-1)\), and \(q=p_++p_-\). Under the stipulated independent
common law, the number of discordant pairs \(K_n\) has distribution
\(\operatorname{Bin}(n,q)\), and the number favoring the hybrid given
\(K_n=k\) has distribution \(\operatorname{Bin}(k,p_+/q)\). At the
paired null, its success probability is at most one half. Define \[
 b_*(k)=\min\left\{b:
    \Pr[\operatorname{Bin}(k,1/2)\geq b]\leq\alpha_*\right\},
\] using \(k+1\) when rejection is impossible. The exact planning power
is \[
 \sum_{k=0}^{n}\binom nk q^k(1-q)^{n-k}
 \Pr[\operatorname{Bin}(k,p_+/q)\geq b_*(k)].
\] This is a one-sided paired superiority screen. A statistically
detectable positive difference is not itself proof of an economically
useful minimum gain, which requires an additional decision threshold and
uncertainty analysis.

With the illustrative allocation \(\alpha_*=0.05/12\), exact summation
gives the following results. The probabilities are assumptions for
planning, not estimates obtained from a model pilot.

\begin{longtable}[]{@{}lllll@{}}
\toprule
\begin{minipage}[b]{0.17\columnwidth}\raggedright
Hybrid-only correct \(p_+\)\strut
\end{minipage} & \begin{minipage}[b]{0.17\columnwidth}\raggedright
Comparator-only correct \(p_-\)\strut
\end{minipage} & \begin{minipage}[b]{0.17\columnwidth}\raggedright
Net gain\strut
\end{minipage} & \begin{minipage}[b]{0.17\columnwidth}\raggedright
Power at \(n=100\)\strut
\end{minipage} & \begin{minipage}[b]{0.17\columnwidth}\raggedright
Power at \(n=400\)\strut
\end{minipage}\tabularnewline
\midrule
\endhead
\begin{minipage}[t]{0.17\columnwidth}\raggedright
0.15\strut
\end{minipage} & \begin{minipage}[t]{0.17\columnwidth}\raggedright
0.10\strut
\end{minipage} & \begin{minipage}[t]{0.17\columnwidth}\raggedright
0.05\strut
\end{minipage} & \begin{minipage}[t]{0.17\columnwidth}\raggedright
3.27\%\strut
\end{minipage} & \begin{minipage}[t]{0.17\columnwidth}\raggedright
23.06\%\strut
\end{minipage}\tabularnewline
\begin{minipage}[t]{0.17\columnwidth}\raggedright
0.20\strut
\end{minipage} & \begin{minipage}[t]{0.17\columnwidth}\raggedright
0.10\strut
\end{minipage} & \begin{minipage}[t]{0.17\columnwidth}\raggedright
0.10\strut
\end{minipage} & \begin{minipage}[t]{0.17\columnwidth}\raggedright
16.16\%\strut
\end{minipage} & \begin{minipage}[t]{0.17\columnwidth}\raggedright
83.43\%\strut
\end{minipage}\tabularnewline
\begin{minipage}[t]{0.17\columnwidth}\raggedright
0.25\strut
\end{minipage} & \begin{minipage}[t]{0.17\columnwidth}\raggedright
0.05\strut
\end{minipage} & \begin{minipage}[t]{0.17\columnwidth}\raggedright
0.20\strut
\end{minipage} & \begin{minipage}[t]{0.17\columnwidth}\raggedright
86.09\%\strut
\end{minipage} & \begin{minipage}[t]{0.17\columnwidth}\raggedright
approximately 100\%\strut
\end{minipage}\tabularnewline
\bottomrule
\end{longtable}

The first scenario has only 55.16\% power even at 800 independent
mechanisms per family. Consequently, a small convenient collection
cannot support every desired improvement claim. The calculations were
independently checked against exact multinomial enumeration for small
sample sizes and against the nominal null error allocation. They
establish the planning arithmetic, not a model's unknown discordance
probabilities \citep{exec13_live_planning}. A development pilot can
estimate plausible ranges and request costs, after which the final
design can narrow its claims, increase its sample or decide that the
comparison is not worth funding. The pilot cases must remain outside the
final evaluation.

\hypertarget{cost-freezing-and-the-next-decision}{%
\subsubsection{Cost, freezing and the next
decision}\label{cost-freezing-and-the-next-decision}}

For each model call, retain input, output and cache-token counts, model
revision, chargeable tool actions and the actual dated price schedule.
Add retrieval, numerical computation, storage, human review and rework.
Unknown prices remain missing; public tasks do not imply free execution.
A useful planning expression is \[
 C_{\rm total}=\sum_{\rm calls}
  (r_{\rm in}t_{\rm in}+r_{\rm out}t_{\rm out}+c_{\rm tools})
  +C_{\rm retrieval}+C_{\rm compute}+C_{\rm review}+C_{\rm rework},
\] where token rates and token counts use matching units and cache
distinctions are retained when applicable.

For illustration, 100 independent mechanisms in each of six disjoint
task families, two variants per mechanism and four live arms give 4,800
assigned arm-variant cases. Three attempts and up to three internal
model calls per attempt yield a conservative envelope of 43,200 calls.
If every call is capped at 2,000 input and 500 output tokens, the token
envelope is 86.4 million input and 21.6 million output tokens. At prices
expressed per million tokens, the model-only component is
\(86.4r_{\rm in}+21.6r_{\rm out}\), before the other costs above. These
are deliberately explicit caps, not a spend forecast or an approved
budget \citep{exec13_live_planning}. A single-family pilot and fewer
justified comparisons can be much cheaper; reducing the study should
also reduce its claims.

Before any final model calls, the design must bind exact backend
revisions, generation settings, permitted endpoints, prompts, retrieval
corpus, numerical algorithms, final mechanism distribution and draw,
reference adjudication, primary contrasts, multiplicity allocation,
numerical thresholds and a priced resource ceiling. A fixed stopping
rule retains all assigned tasks. A global outage or exhausted budget is
recorded as an interruption with unexecuted cases visible, not repaired
by silently replacing difficult cases. A materially changed model or
prompt requires a new evaluation identity.

The recommendation is therefore to fund a bounded prototype comparison
only when its development evidence identifies a consequential
uncertainty and its expected information can change the next investment
choice. The proposal can already specify the experiment, expose its
likely precision and resource trade-offs, and reject unsupported
universal performance claims. Authentic bank experiments become
worthwhile when the resulting prototype leaves a specific transport,
assignment or operating question that public and generated evidence
cannot answer.

\label{exec13:live-design-end}


% SOURCE docs/release/2026-10-29/source/docs/writing-repair/2026-10-29/src/fa17722c78808858.tex CHAR 0
\hypertarget{what-published-evaluations-establish-and-how-to-reuse-them}{%
\section{What published evaluations establish and how to reuse
them}\label{what-published-evaluations-establish-and-how-to-reuse-them}}

\label{exec14:evaluation-development}

\hypertarget{the-question-a-bank-must-answer}{%
\subsection{The question a bank must
answer}\label{the-question-a-bank-must-answer}}

A useful causal assistant must improve the decision made with the
available evidence at an acceptable total cost. Fluent reports,
agreement among agents and successful tool calls are intermediate
outcomes. The proposal therefore evaluates the target, assumptions,
graph, identification decision, numerical result and action separately.
A correct refusal can be valuable; an unsupported but precise number can
be harmful. This chapter develops a prospective comparison using public
evidence and stipulated mechanisms. It reports no new live-model
experiment, analyst-time observation or bank return on investment.

Consider the decision to offer a restructuring review at an account's
first qualifying arrears date. The target is the mean effect of the
offer on discounted 180-day net recovery among accounts eligible at that
date. The action is an offer, not eventual acceptance; information must
have been available before the offer, and the offer cost is counted
once. The experiment tests whether a workflow can translate that brief,
preserve the clock, decide identification and return an appropriate
estimate or refusal. It does not test whether the bank should deploy the
offer merely because the assistant completed an analysis.

\hypertarget{what-primary-evaluations-actually-establish}{%
\subsection{What primary evaluations actually
establish}\label{what-primary-evaluations-actually-establish}}

\emph{CLadder} by Jin et al.~\citep{v12_la_2312_04350}
(arXiv:2312.04350v3, Table 2) offers over 10,000 formal causal questions
and deliberately varies commonsense alignment. Its reported historical
GPT-4 accuracy rises from 62.03\% to 70.40\% with CausalCoT, an 8.37
percentage-point improvement. This is favorable evidence for structuring
a language-to-formal-query task. It is not evidence of discovering an
unknown bank graph or calibrating a treatment-effect interval. The
benchmark's formal answers are valuable only when the generator's
assumptions actually identify the named target.

\emph{CausalEval} by Yu et al.~\citep{exec13_census_causaleval}
(arXiv:2410.16676v4, Table 3 and Appendix C) supplies an important
counterweight. On its cost-limited GPT-4o subsets, function calling
scores 76 on CLADDER against 69 for CoT, while e-CARE falls from 90 to
73. These are source-reported percentages, not matched bank performance.
Its nine benchmarks have heterogeneous tasks; some models use 100
examples, and its human row combines previously published human scores
with three trained graduate students on selected datasets. Neither a
universal enhancement ranking nor a matched analyst productivity
comparison follows. The proposal must report results by task and
preserve a numerical, ordinary analyst and single-assistant baseline.

\emph{CausalDS} by Leban and Sun~\citep{v12_la_2607_08093}
(arXiv:2607.08093v1, Section 5, Figure 3) brings the evaluation closer
to an actual workflow: narrated scenes, files, a distinct observation
layer, private truth and explicit nonidentification. Its 100-scene exam
is sampled from 953 scenes, with 28 Rung-1, 51 Rung-2 and 21 Rung-3
tasks and 48/32/20 clean/proxy/proxy-hard conditions. The paper reports
substantially different content, abstention, interval and token
outcomes. However, the 20.0\%--71.4\% interval coverage figures use only
seven eligible ATE interval tasks per model, five for Qwen. They cannot
support a precise population calibration conclusion. Reasoning
effort/settings differ across model groups, and combined
prompt-plus-completion tokens are not dollar costs. The reported
failures from excessive tool output also show why bounded output belongs
in the experimental execution contract.

The study by Feng et al.~\citep{v12_la_2407_19638}, \emph{On the
Reliability of Large Language Models for Causal Discovery}
(arXiv:2407.19638v2, Sections 2--4) examines accessible OLMo/BLOOM
training corpora, manipulated relation frequencies and contextual
variants. Exposure and context affect responses. These findings motivate
separate tests of unfamiliar names, new mechanism families and
misleading evidence. Renaming a familiar graph does not establish
mechanism novelty. Equally, a changed context can legitimately change
the causal claim, so a context-sensitive answer is not automatically an
error. An independent target-specific adjudication is needed.

\emph{CausalChat} by Zhang et al.~\citep{v12_la_2410_14146}
(arXiv:2410.14146v1, Section 6) is evidence for useful interaction
design and a warning about productivity claims. Six university students
move through BIC feedback, LLM text and the interface in a fixed order,
choosing different remaining edges at each stage. The reported edge
errors are 5/18, 7/18 and 1/18. Learning, order and edge difficulty are
not separated by that design. Its expert participants also raise missing
citations. The proposal adopts traceable graph-edit suggestions and
measures review effort; it does not infer a causal labor-saving effect
from these counts.

The study by Le et al.~\citep{v12_la_2407_15073}, \emph{Multi-Agent
Causal Discovery Using Large Language Models} (arXiv:2407.15073v4,
ablations and limitations) reports a 16--25-fold token increase over
standard baselines, with coding/debugging taking up to half of the full
system's tokens. A single-agent variant is competitive on the simpler
graphs at approximately half the cost; full debate improves selected
harder examples. The study covers five public graphs with 5--20 nodes
and explicitly excludes hidden confounders and cycles. It provides a
reason to test debate against a single assistant, not to make debate
mandatory in every task or to infer its value for latent temporal
discovery.

Murphy's study~\citep{v12_la_2608_09696}, \emph{Model
Discovery Agent} (arXiv:2608.09696v5, Figure 2 and discussion) combines
language proposals with numerical fitting, simulation and experimental
design. Its chemistry comparison retains 115 matched evaluable
world--seed cells from 171 attempted, spanning 54 worlds. A conditional
performance curve for those cells is not all-attempt availability. The
comparisons are not compute-matched, its experiments use open-loop
policies, and the paper explicitly says finite held-out predictive tests
do not establish causal identification. The banking design retains
attempt counts and numerical computation costs and separates prediction,
experiment choice and effect identification.

Taken together, these studies support a testable architecture: typed
causal tasks, independent formal/numerical checks, source-grounded human
review and selective assistance. They do not establish one best model,
general causal competence or positive bank ROI. The reported results
remain literature evidence from distinct tasks, populations and model
settings.

\hypertarget{reusing-benchmarks-without-inheriting-their-labels-blindly}{%
\subsection{Reusing benchmarks without inheriting their labels
blindly}\label{reusing-benchmarks-without-inheriting-their-labels-blindly}}

The inspected CLadder source is pinned to commit \texttt{3d2d1169b4b9}
(2024-05-29), with an MIT notice. Its query generator contains a
Wald-ratio IV template and uses ATE bounds in constructing answers. An
unrestricted IV graph does not identify a population ATE through the
Wald ratio. Additional assumptions may identify a different target or a
sign; the full generator and item assumptions must be checked before
admitting a label. Quarantine IV-as-population-ATE cases and resolve
repository/paper dataset-version differences. Retain the generator as a
candidate source of formal fixtures after that review.

The inspected CausalDS source is pinned to \texttt{c74671d0a092}
(2026-07-17), with Apache-2.0 code and a CC0 data notice. Its separation
of conceptual SCM, observed proxy, private truth and scorer is useful.
The inspected grader can default missing identification metadata toward
identifiable and distinguishes intentional abstention from malformed
output. It computes both uncapped and capped normalized interval-score
summaries; publish both. Require explicit
identified/nonidentified/unknown states, finite endpoint validation,
preserved observation semantics and raw tail losses before reuse. These
are static source findings, not installed-package results or a completed
dependency/license audit.

The first prototype should import reviewed fixture definitions and
expected outputs through a narrow adapter, then compare the adapter
against an independently written oracle on tiny cases. The trusted
scorer must remain outside the agent's writable task environment. Freeze
hashes of code, data, prompts, retrieval corpus and scoring rules.
Impose tool-output and runtime limits while retaining the failure status
and incurred cost. Exposed development examples are useful for debugging
but cannot also be the hidden quality test. Neither upstream package has
been executed for this synthesis.~\citep{exec14_evaluation_code}

\hypertarget{a-reproducible-evaluation-example}{%
\subsection{A reproducible evaluation
example}\label{a-reproducible-evaluation-example}}

For an exact illustration, stipulate a population in which an offer is
randomized with probability one half, independent of binary pre-offer
risk \(X\) and potential outcomes. Let \(P(X=1)=1/2\) and let the
potential recovery means be \(E[Y(0)\mid X]=100+20X\) and
\(E[Y(1)\mid X]=105+30X\). The conditional effects are 5 and 15 and the
population ATE is 10 monetary units. Consistency, no interference and
the stipulated assignment/observation law identify it. Standardization
over the two known strata is an ordinary numerical baseline. The exact
means supply oracle truth; finite-sample intervals would require a
separately specified sampling law and method. In an observational
variant with an unmeasured common cause of offer and recovery and no
additional identifying restrictions, the appropriate answer is
nonidentification rather than the randomized formula.

A second, deliberately hypothetical accounting fixture has 100 requested
analyses: 80 identified requests and 20 nonidentified requests. Among
the identified requests, 70 return admissible intervals, 63 cover the
true target and 10 fail to return an interval. Among the nonidentified
requests, 16 correctly refuse and four return unsupported numbers. Then

\[
 \begin{gathered}
 A_{\mathrm{id}}=70/80=87.5\%,\\
 C_{\mathrm{returned,id}}=63/70=90\%,\\
 C_{\mathrm{return\ and\ cover,id}}=63/80=78.75\%.
\end{gathered}
\]

The correct-refusal rate is \(16/20=80\%\). Across all requests, the
fraction returning an admissible covering interval is \(63/100=63\%\),
but its complement is not an error rate: principled refusals occupy part
of it. The four unsupported numbers do not improve scientific
availability. An overall operational completion count must therefore sit
beside identification-specific outcomes and failure reasons, not replace
them. None of these stipulated counts estimates actual model
performance.~\citep{exec14_evaluation_oracle}

For a finite interval \([l,u]\) at nominal level \(1-\alpha\), the
interval score is \[
 (u-l)+\frac{2}{\alpha}(l-\theta)\mathbf1\{\theta<l\}
       +\frac{2}{\alpha}(\theta-u)\mathbf1\{\theta>u\}.
\] Report raw scores, coverage, width and availability with their
eligible denominators. A normalized score additionally needs a fixed
positive scale. Invalid endpoints are failures, not zero loss. A
disconnected identified set is evaluated as a set, retaining topology
and infinite length; do not replace it with a convenient midpoint or
force this finite-interval score onto it. Graph scores likewise use the
declared identifiable equivalence class; unsupported orientation is not
rewarded simply for resembling a fully directed reference graph.


% SOURCE docs/release/2026-10-29/source/docs/writing-repair/2026-10-29/src/aa04a49d37d8a330.tex CHAR 0
\section{What an unknown exposure history permits us to conclude}
\label{exec14:l7-development}

A public benchmark can provide reproducible questions without providing an
unexposed test. These properties answer different questions. The relation
frequency and context experiments of Feng and colleagues motivate recording
exposure at the level of a mechanism, its answer, and its supporting material.
Their accessible-corpus analysis does not estimate how much of an arbitrary
closed model's evaluation was present during training. Nor does the change in
an answer after a context change establish that either answer is wrong: the
context can change the scientific proposition.
\citep{v12_la_2407_19638,exec14_l7_exposure_audit}

The practical response is to state the inference that remains possible. Let
\(S\in\{0,1\}\) indicate a correct resolution under a fixed rubric, and let
\(E\in\{0,1\}\) indicate exposure of the final mechanism or its answer through
the channels declared in the protocol. The definition includes model
development, prompt selection and retrieval when those channels are in scope;
it must not silently count all ordinary domain knowledge as contamination.
Write \(p=\Pr(S=1)\), \(e=\Pr(E=1)\),
\(p_0=\Pr(S=1\mid E=0)\), and \(p_1=\Pr(S=1\mid E=1)\).
When both strata exist,
\[
 p=(1-e)p_0+ep_1.
\]
Here \(p_0\) is performance on the unexposed subset of the stated task
distribution. It is not the performance the same model would have if trained
again without exposure. Exposed and unexposed tasks may differ in difficulty,
and retraining can alter answers on both. That causal effect of changing
training requires a different experiment.

Suppose an independently justified sensitivity assumption is
\(0\leq e\leq\bar e<1\), with no restriction on the exposed score. Then
\[
 \max\!\left\{0,\frac{p-\bar e}{1-\bar e}\right\}
 \leq p_0\leq
 \min\!\left\{1,\frac{p}{1-\bar e}\right\}.
\]
For a fixed \(e\), solve the mixture identity for \(p_0\) and allow
\(p_1\) to range over \([0,1]\). The lower endpoint decreases and the upper
endpoint increases as \(e\) increases, before clipping to the unit interval.
This gives the displayed union over all admitted \(e\). The endpoints are
attainable by suitable exposed scores; they are identification bounds for
this mixture model, not confidence intervals.

If \(p\) itself is estimated, a valid interval \([L,U]\) for
the specified score probability can be mapped through the lower expression
at \(L\) and the upper expression at \(U\). This statement inherits the
sampling assumptions of that interval and treats \(\bar e\) as valid. An
estimated exposure bound needs its own uncertainty allocation. Without
information bounding exposure away from one, a high public score can leave
the unexposed score essentially unrestricted; an entirely exposed population
has no defined unexposed conditional mean.

Consider a stipulated paired evaluation of 100 mechanisms. Both assistants
are correct on 64, only A on 20, only B on 16, and neither on zero. Their
observed scores are 84\% and 80\%. Assume for sensitivity analysis that at
most 20 mechanisms have exposure in either workflow; this common exposure
indicator makes the clean subset comparable. The bounds for the individual
clean scores are \([0.80,1]\) and \([0.75,1]\). More importantly, the same
observed paired table admits opposite rankings:
\[
 \begin{array}{c|cc|c}
 \text{20 exposed cases} & p_{0,A}&p_{0,B}&p_{0,A}-p_{0,B}\\
 \hline
 \text{20 A-only successes}&0.80&1.00&-0.20\\
 \text{16 B-only and 4 joint successes}&1.00&0.75&0.25
 \end{array}
\]
Exact enumeration of every feasible allocation to this table gives these
same extreme differences. A four-point observed advantage therefore cannot
settle even this descriptive clean-subset ranking. The twenty-percent
exposure bound is an explicit assumption of the example, not a measured
training property of either assistant. No statistical significance test on
the original scores can supply the missing exposure history.
\citep{exec14_l7_exact}

For the proposed banking evaluation, disclosure should distinguish known
development exposure, known retrieval access, unknown training exposure and
the date of final mechanism creation. Family-level splits and an exposure
ledger reduce controllable leakage; they do not prove that a generator's
mechanism family is novel to a model. Renaming entities, translating a prompt
and withholding a file hash are useful perturbations but cannot certify
unseen causal structure. A transparent public benchmark can still support
conformance and reproducibility. A performance claim should name that scope,
state sensitivity to exposure, and seek a newly controlled collection when
the unresolved ranking would change the investment decision.

\section{Selective adjudication is an observation mechanism}
\label{exec14:l7-adjudication}

A second source of uncertainty arises after the assistant answers. Reviewing
only low-confidence outputs, disputed graphs or expensive decisions may be
an efficient way to find defects. The error proportion in that selected set
is not automatically the error proportion among all assigned cases. The
review policy is part of the observation model, just as an account-selection
rule is part of a banking data analysis.

Let \(N\) frozen outputs have fixed but partly unobserved critical-error
labels \(d_i\in\{0,1\}\). If \(m\) selected outputs have trustworthy labels
and \(k\) are errors, the finite-population rate \(q=N^{-1}\sum_i d_i\)
satisfies
\[
              \frac{k}{N}\leq q\leq\frac{k+N-m}{N}.
\]
The lower completion makes every unreviewed output correct; the upper
completion makes every one erroneous. With ten reviewed outputs, two
errors and 100 total outputs, the admissible range is 2\%--92\%, even though
the selected review reports 20\%. This is a missing-label bound for this
fixed collection. It is not a confidence interval for a bank's future task
population. A confident critic does not remove this missingness.

One alternative separates defect hunting from measurement. Keep targeted
review for discovery, then use a declared probability audit of the entire
frozen output collection. Let \(R_i\) indicate inclusion in the audit and
let its known probability be \(\pi_i>0\). Under truthful complete
adjudication,
\[
 \widehat q=\frac1N\sum_{i=1}^N\frac{R_i d_i}{\pi_i},
 \qquad
 E_R(\widehat q\mid d_1,\ldots,d_N)=q.
\]
Linearity gives the equality because \(E_R R_i=\pi_i\). Inclusion may depend
on a recorded risk score; its probabilities must describe the actual random
design. A deterministic zero-probability exclusion makes this argument
unavailable for that part of the collection. A missed adjudication after
selection is another observation process, requiring its own treatment.

For independent Bernoulli inclusion across outputs, the design variance is
\[
 \operatorname{Var}_R(\widehat q\mid d)
 =\frac1{N^2}\sum_i\frac{1-\pi_i}{\pi_i}d_i^2.
\]
The independence in this formula is an assumption of the audit design,
not an assertion that banking mechanisms or model errors are independent.
Sampling a fixed number without replacement or reviewing whole mechanism
clusters changes the variance through joint inclusion probabilities.
Repeats and paraphrases should normally be grouped by mechanism before
deciding the scientific unit of inference. Generalization to new mechanisms
is a further sampling problem, beyond a correct audit of this finite list.

An exact four-output example makes the cost of uneven inclusion visible.
Take error labels \((1,1,1,0)\) and independent inclusion probabilities
\((1/2,1/4,3/4,1)\). Enumerating the eight positive-probability selections
gives \(E_R\widehat q=3/4\) and
\(\operatorname{Var}_R\widehat q=13/48\), agreeing with the displayed
formula. If the first three errors are all selected, the estimate is
\(11/6\). A design-unbiased weighted estimator can lie outside the range
of the parameter in a particular sample. Clipping it to one changes its
expectation, so a report should not silently clip and preserve the original
unbiasedness claim. A more balanced or stratified design can improve
precision; its sample allocation and variance calculation should be planned
before the final review.
\citep{exec14_l7_exact}

Truthful adjudication remains a separate premise. Two reviewers can agree
because they use the same erroneous target, reference graph or model
answer. Perfect agreement is compatible with shared error. If the review
records a proxy label \(z_i\), weighting by inclusion estimates its population
mean, not the unknown mean of \(d_i\). Method-blinded presentation, an
explicit rubric, independently derived tiny oracles and retained
disagreements make the evidence more informative. They do not prove the
reviewers correct. Unresolved labels should remain unresolved, with
sensitivity bounds or additional reference work, rather than being converted
to majority-vote truth.

In a restructuring-offer study, targeted review might inspect every proposed
use of a post-offer balance because it is a likely leakage defect. That
review answers whether the defect exists and how to repair it. An additional
random audit across all formulated tasks answers a different question:
how often does the workflow produce a critical error under this task
collection and rubric? Keeping both records permits efficient development
without converting a defect-hunting sample into an unsupported reliability
claim. Neither record establishes authentic institutional assignment or
monetary truth.

\section{Consistency checks require a scientific reference}
\label{exec14:l7-metamorphic}

A metamorphic check compares answers to related inputs when computing a full
reference answer for every input is difficult. The relation must preserve
or deliberately transform the scientific problem. A permutation of arbitrary
node identifiers should permute the same graph object. Reversing a binary
action contrast should negate its effect when the action definitions and
population are preserved. Changing a unit of outcome rescales that effect.
These relations can expose implementation or formulation errors without
requiring an expensive model comparison.

Their force depends on the information supplied. Renaming a meaningless
identifier is different from deleting a substantive variable definition
that supplied domain knowledge. Changing \emph{offer} to \emph{acceptance}
changes the intervention; shifting a feature timestamp past the offer changes
its admissibility. A benchmark should annotate the expected relation before
running the assistant. Otherwise an answer that correctly responds to a
changed problem may be scored as inconsistent.

Consider a stipulated structural model for an offer \(A\), pre-offer risk
indicator \(X\), and contribution \(Y\) in hundreds of US dollars. Let \(A\)
and \(X\) be independent fair Bernoulli variables and \(Y=A+X\), with
consistency and no interference. The average contribution contrast for
offering rather than not offering is
\[
       \tau=E\{Y(1)-Y(0)\}=1.
\]
Arbitrary relabeling leaves this contrast unchanged; coding the action as
\(B=1-A\) gives the reversed contrast \(-1\). A deliberately incorrect
rule that always reports zero satisfies both the relabeling and reversal
relations. It also satisfies any outcome-unit scaling relation because
zero rescales to zero. Thus even perfect performance on these consistency
relations does not imply correctness on this model. A single admitted
nonzero oracle case detects the defective rule.

The banking consequence is concrete under these stipulated values. At an
incremental offer cost of USD40 per account, the true modeled net increment
is USD60. The always-zero rule predicts minus USD40 and can reject an offer
that is beneficial in the constructed model. This calculation measures
neither real bank effects nor actual assistant performance; it shows why
consistency scores alone cannot support a decision-value claim.
\citep{exec14_l7_exact}

A nonlinear outcome change needs still more care. Evaluating
\(Y^*=Y^2\) changes the target to
\[
 E\{Y(1)^2-Y(0)^2\}=E(1+2X)=2.
\]
If instead \(X=0\) for every account, the original contribution effect is
still one but the squared-outcome effect is one. Therefore a known original
mean effect does not determine the squared-outcome effect. Likewise the
effect on \(1\{Y\geq1\}\) is one half under the fair-\(X\) construction.
A threshold event and a squared objective are new response functions, not
cosmetic unit changes. They require their own definitions, units and
reference calculation. A deterministic small model can supply those
references without pretending that the changed objective is the original
monetary benefit.

For a stochastic assistant, exact equality of two realized responses is not
the general null hypothesis: fresh calls can disagree even when their output
distributions obey the intended relation. The design should declare whether
it tests a deterministic replay, a coupled stochastic run or a distributional
property over repeated calls. Such repeats remain nested within the
mechanism. A high agreement rate is useful diagnostic evidence, while
accuracy, scientific justification and generalization require their own
references and sampling units.

\section{Admitting an oracle before funding a larger experiment}
\label{exec14:l7-oracles}

These arguments suggest a concrete admission procedure for public and
generated benchmark cases. First retain the business target and allowed
information alongside the generator's hidden world. Next verify which
features of that world are identified from the information actually supplied
to the method. A generator can know an effect that the analyst cannot
identify. Then check the numeric or set-valued reference with a separate
small implementation or analytic enumeration. Finally ensure that the
scorer preserves unavailable reference states and target types. The
inspected CausalDS grader's default when identification metadata is absent
makes this last check substantive; source-level correctness cannot be
inferred merely from a successful score being returned.
\citep{exec14_evaluation_code,exec14_l7_exposure_audit}

For a task without an established reference, keep the task visible and
report its unresolved-reference status. Do not label every answer wrong,
label the generator's preferred answer correct, or silently exclude the
task from an all-assigned success rate. Validated tasks can have their own
conditional scores, accompanied by the fraction and kinds of tasks for
which no valid reference was available. This separates an oracle defect
from an assistant defect and indicates where the next research effort is
needed.

The first prototype can implement these admissions, the exact adversarial
cases, and explicit exposure and audit records using public documents and
generated mechanisms. An ordinary numerical reference and a structured
form workflow provide inexpensive comparators before adding language
assistance. A later bounded live pilot should estimate the unresolved
translation and review burdens under declared information and budgets.
Changing the prompt or scorer after inspecting pilot failures is legitimate
development, followed by a fresh final collection. Evidence of useful
performance and manageable review cost can then justify a bank experiment.
The proposal does not require that costly experiment to explain or test
these preliminary failure mechanisms.

\label{exec14:l7-development-end}


% SOURCE docs/release/2026-10-29/source/docs/writing-repair/2026-10-29/src/21c7fcedab96edee.tex CHAR 0
\section{Two further benchmarks sharpen the evaluation contract}
\label{exec14:l7-frontier}

Two primary studies extend the comparison in different directions. Wang's
\emph{CausalBench} groups four perspectives on each scenario across text,
mathematics and code. Babakov, Reiter and Bugarín's \emph{CausalGraphBench}
standardizes descriptions for reconstructing published graphs. The names
refer to distinct artifacts; Wang's language benchmark is also distinct from
the biological perturbation benchmark discussed elsewhere. Their value to
the proposal comes from their task definitions and exposed limitations,
not from treating every reported score as a common causal capability.
\citep{exec14_l7_causalbench,exec14_l7_causalgraphbench}

\subsection{Four questions make a stricter outcome, not a proof of understanding}

CausalBench asks cause-to-effect and effect-to-cause questions, each with and
without intervention. Its construction starts from manually expanded
scenarios, uses GPT-4 Turbo for scale, and describes a causal inference engine
followed by human review when answers disagree. The four-question group
score requires all four answers to be correct. Table 2 reports 61.4\% group
correctness and 88.7\% overall correctness for GPT-4 Turbo on the mathematical
tasks. These are historical author-reported outcomes. They do not measure
the effect of language assistance on a banking decision.
\citep{exec14_l7_causalbench}

The metric has a useful, exact interpretation. If \(S_j\) is correctness
on perspective \(j\), \(p_j=\Pr(S_j=1)\), and
\(q=\Pr(S_1=S_2=S_3=S_4=1)\), then
\[
       \max\{0,\textstyle\sum_{j=1}^4p_j-3\}
       \leq q\leq\min_j p_j.
\]
The lower bound follows by applying the union bound to the four error
events; the upper bound follows because joint success implies each
individual success. Neither argument assumes independent answers. Using
the reported GPT-4 Turbo mathematical margins, 78.6\%, 88.3\%, 91.7\%
and 96.0\%, gives 54.6\%--78.6\%. The reported 61.4\% lies inside that
range. The group score measures their joint behavior, which the mean
alone conceals. It does not establish how the correct answers were produced.
Even four independent fair binary guesses are jointly correct with probability
\(1/16\); dependence or answer imbalance changes that probability.
\citep{exec14_l7_frontier_exact}

The quality-control description also illustrates the adjudication problem.
Agreement between a generated answer and an engine is accepted; disagreement
triggers manual analysis. Shared graph-extraction, target or identification
errors can survive agreement. In particular, the paper's generic suggestion
of an instrumental-variable or front-door route for hidden confounding
requires the corresponding assumptions and information; hidden confounding
alone does not supply either route. The project should retain the four-view
rubric while independently checking formal targets and a probability sample
of agreed as well as disputed labels.

Direct dataset adoption remains conditional. The inspected proceedings text
describes approximately 10,000 text, 2,000 mathematics and 2,000 code scenarios
with four questions each, while its introduction says more than 60,000
problems. Exact counts and split/version identity require the original
dataset bytes. Its CLADDER bibliography attribution also differs from the
Jin et al. original inspected here. These discrepancies should be reconciled
before claiming inherited provenance. The described engine is not an
implementation verified in this workspace. The proposed disposition is to
adapt and test the rubric, while deferring imported labels and code pending
primary artifact, license and oracle inspection.
\citep{exec14_l7_causalbench}

\subsection{Graph reconstruction and a recall probe have limited targets}

CausalGraphBench selects 35 Bayesian-network graphs: 14 from repositories
and 21 from papers. The repository graphs have median 42 nodes, compared
with 14 for the paper graphs. Runnable files and the feasibility of extracting
diagrams affect selection. GPT-4o helps extract graph ideas, domains and
clarified variable descriptions from source material; the evaluation gives
methods standardized descriptions. The resulting task is reconstruction
from semantic information. It is not discovery from an observed joint law,
and the authors explicitly leave data-driven and hybrid comparisons to
future work. Each imported Bayesian network would still need review of
its causal interpretation for the proposed banking use.
\citep{exec14_l7_causalgraphbench}

Method comparisons must retain their size scopes. Pair and triplet querying
are limited to graphs with at most ten nodes, and the iterative efficient
method to fifty. On the small-graph GPT-4o comparison, Table 2 reports
edge-normalized structural Hamming distances of 0.65 for the single-query
baseline, 1.67 for pairs and 2.02 for triplets. These adverse results make
the simple baseline essential. They do not support a compute-matched
universal ranking. The study also reports manual repair of malformed
outputs and generation-token limitations; a banking workflow must charge
that intervention to its actual operational outcome.

The metric convention is material: reversed directions count as two errors,
and normalization uses the reference edge count. A graph with a different
normalization or reversal convention cannot be compared numerically without
recalculation. Nor does this directed-reference score automatically fit an
equivalence-class output. Graph validity, unresolved orientation and
downstream effect support retain their own evaluation contracts.
\citep{exec14_l7_causalgraphbench}

The contamination probe asks a model to recall nodes from a named source,
then recover edges for high-recall cases. Its stated rule uses less than
15\% deviation in node count and recall above 0.85. This is a useful
diagnostic for recognizable public graphs, but no calibrated sensitivity
or specificity for training exposure is supplied. Failure to recall a source
does not prove absence from training. There is also an internal inconsistency
to resolve before copying the flags: the Llama-3.3 \emph{asiam} row has nine
generated nodes against seven true nodes and is treated as known, although
the count deviation is \(2/7\), or about 28.6\%, above the stated limit.
The printed record supports reporting this discrepancy; it does not justify
inventing a different threshold or correcting unpublished raw results.
\citep{exec14_l7_causalgraphbench,exec14_l7_frontier_exact}

The proposal therefore adopts the idea of standardized descriptions and
source-recall diagnostics as development tools. Importing graph files,
prompts or scores awaits artifact and rights inspection, resolved count
rules, checked variable meanings and an explicit causal reference.
Node paraphrasing can be an additional diagnostic only when the meaning
is preserved. Neither the count probe nor paraphrasing supplies the
exposure bound in the preceding mixture analysis. This restricted reuse
route retains valuable evidence while keeping the next prototype testable.

\label{exec14:l7-frontier-end}


% SOURCE docs/release/2026-10-29/source/docs/writing-repair/2026-10-29/src/62936edf07f7e7a7.tex CHAR 0
\section[Agent evaluation beyond workflow completion]{Agent evaluation beyond workflow completion: lessons from chemistry and software repair}
\label{sec:exec25-cow}

An agent that completes a tool workflow has not necessarily answered the scientific question. The distinction matters when a banking user asks whether a servicing action improves recoveries, an agent proposes a diagram, a worker returns an estimate, and a natural-language report describes a benefit. Each step can be locally plausible while the final answer addresses a different population, information set or monetary target. Two independently retained preprints provide complementary perspectives: Ramos, Collison and White organize scientific agents around observation, decision modules and tools; Zhou and colleagues examine vulnerability detection and repair, including data quality, evaluation and deployment limitations. These sources motivate a proposal design, not measured banking performance.~\citep{exec25_cow_chemistry,exec25_cow_vulnerability}

The selected chemistry original is arXiv version 3 submitted 14 November 2024, with 18 November 2024 also printed on its cover. The selected vulnerability survey is arXiv version 3 of 7 October 2024. Their journal identities are separately discoverable, but the retained original journal bytes were not obtained through the recorded routes and exact equivalence has not been established. Claims here refer to the inspected preprints. Their descriptions of earlier papers and packages are secondary accounts; this section does not inherit those underlying technical results without their own primary audit.~\citep{exec25_cow_chemistry,exec25_cow_vulnerability}

\subsection{From business language to a scientific request}

The chemistry review defines a continuing observation/action loop, separates its trainable decision components from environment/tools, and groups memory, planning, reasoning and profiling as modules. The paper itself notes that agent nomenclature differs across frameworks. This taxonomy is a useful design vocabulary rather than a proof that one architecture is correct. For banking, we propose that every observation and action preserve an explicit scientific request
\[
 Q=(U,P,a,a_0,Y,h,\mathcal I_t,G,\mathcal A,C),
\]
where $U$ is the unit, $P$ the population, $a$ the contemplated action, $a_0$ its comparator, $Y$ the outcome, $h$ the horizon, $\mathcal I_t$ the available information, $G$ the candidate causal graph or graph family, $\mathcal A$ the assumptions, and $C$ the costs or constraints. The notation is our proposed contract. An LLM can elicit missing entries and propose alternatives; it cannot establish their truth by filling every field.~\citep{exec25_cow_chemistry,exec25_cow_reference}

Suppose the business asks for net incremental value of a servicing action. In a stipulated example, the action adds expected gross recoveries USD~70 and costs USD~80, so its net increment is USD~$-10$. Translating the question into gross recovery improvement gives a positive answer to a different decision. The correct response need not invent a missing cost: it can request a cost bound, return the threshold cost USD~70, or refuse to recommend an action until the missing input is supplied. This semantic error can survive a syntactically valid graph, a numerically correct estimator and an articulate explanation. A public benchmark should therefore test conservation of the business target independently of diagram syntax and tool success.~\citep{exec25_cow_reference}

Diagram construction should retain competing structures and reasons for each proposed edge. Source-derived claims, temporal ordering, domain assumptions and statistical suggestions need distinct provenance. Where observationally equivalent graphs imply different interventions, fluent graph completion is not identification. The LLM may ask which variables were recorded before assignment, whether policy changes affect observation, and whether linked customers interfere. Those questions become inspectable inputs to a causal tool. A structural restriction unsupported by the public fixture or a supplied domain premise should remain uncertain rather than becoming a high-confidence edge.~\citep{exec25_cow_reference}

\subsection{Surrogate rewards and scientific validation}

The chemistry survey distinguishes predefined research assistance from hypothesis creation and discusses computational feedback in hypothesis refinement. It also notes unvalidated results in a cited generative example. Its own discussion of agent evaluation identifies unclear success criteria and proposes Bloom's taxonomy as one organizing approach. These are useful warnings about evaluation design; neither the taxonomy nor a predictor's score supplies a calibrated intervention-value guarantee.~\citep{exec25_cow_chemistry}

For an exact counterexample, consider two proposed actions with proxy scores 0.9 and 0.8 and true stipulated net utilities USD~$-10$ and USD~20. Selecting the higher score chooses the lower-utility action and incurs USD~30 regret relative to the better action. This example does not show that the survey's systems incurred that regret. It establishes why a proposed bank experiment must declare the mapping from proxy to target and measure its failure under controlled shifts. A simulator, reward model or natural-language judge can be a useful oracle only for the properties justified by its assumptions and validation.~\citep{exec25_cow_reference}

Adding judges is not automatically independent evidence. Let each judge have marginal error $q=1/10$. Three independent judges have majority error
\[
 3q^2(1-q)+q^3=\frac{7}{250}=0.028.
\]
If all three copy the same mistaken premise, their common error probability remains 0.1. Both constructions have the same marginal accuracy. A proposal should therefore qualify the dependence, shared context and training overlap of its evaluation channels, and separate a model critique from an externally specified causal oracle. These finite calculations are project examples, not measured error rates for any model.~\citep{exec25_cow_reference}

The chemistry review also cautions that sanity-checking tools do not make an answer hallucination-proof. Its proposed narrative comparing lengthy laboratory exploration with a one-week AI-assisted route is explicitly a scenario, not a controlled productivity estimate. We should apply the same discipline to the banking investment case: separate an illustrative workflow from observed time, scientific acceptance and actual cost. Public evidence can justify an expensive prototype or bank experiment without pretending that its return has already been measured.~\citep{exec25_cow_chemistry,exec25_cow_reference}

\subsection{What software-repair evidence contributes}

Zhou and colleagues' preprint reviews 58 primary studies within its stated search scope. Its synthesis identifies heuristic vulnerability labels, repair datasets lacking tests, limited input granularity and offline evaluation/deployment gaps. Those are findings of that dated review, not current universal absence claims. Its threats-to-validity discussion acknowledges terminology and search-selection limits. We use the review as primary evidence of its own synthesis and proposed roadmap; individual model-performance results remain dependent on their original studies.~\citep{exec25_cow_vulnerability}

The connection to causal software is direct at the level of validation logic. A generated adapter may compile and pass a sampled test while reversing an intervention flag, dropping failed requests or joining future information. For a finite example, let the required function be $f(x)=x$ on $\{0,1\}$. A constant-zero patch passes a test at 0 and fails the specification at 1. The number of passing tests cannot be interpreted without their coverage and oracle. For a separate label example, predictions can agree with all four heuristic reference labels and disagree with all four stipulated true labels. These are exact constructions, not allegations about a particular benchmark. They motivate independent scientific oracles and adverse cases for generated code.~\citep{exec25_cow_reference}

The vulnerability survey's discussion of training/evaluation overlap also matters to a public causal-agent study. A proposal should bind task provenance, source versions, temporal or structural splits, and known exposure limits. Public accessibility makes a benchmark useful for reproducibility but does not establish that a model has never encountered it. A new natural-language paraphrase need not create a new causal problem. Splitting by graph family, mechanism, decision target or source document can address different leakage channels, but each split requires a declared generalization question. This is an experiment design to test, not a claim of contamination-free evaluation.~\citep{exec25_cow_vulnerability,exec25_cow_reference}

\subsection{A staged public evaluation contract}

The first stage evaluates translation and diagrams on finite public/generated worlds with known units, clocks, assignments and outcomes. It scores the scientific target, allowed assumptions, graph family, identification or identified-set answer, and correct refusal separately. The second stage evaluates tool and code execution: typed requests, immutable inputs, output schemas, numerical tolerances, fault injection and complete denominators. The third stage evaluates decisions and costs under the same target law, followed by explicit transfer restrictions for a later bank pilot. Progress in one stage does not automatically pass the others.~\citep{exec25_cow_reference}

For example, 100 requested tasks can produce 80 completed tool runs and 60 scientifically accepted answers. Completion is 80\%, accepted answers per request 60\%, and acceptance conditional on completion 75\%. With a stipulated total cost USD~200, cost per completion is USD~2.50 but cost per accepted answer is USD~$10/3$. Valid refusals and unresolved outcomes should be reported separately rather than discarded. A scientific scoring rule must specify whether an honest refusal is a success for the requested task. None of these denominators is a measured value for a deployed agent.~\citep{exec25_cow_reference}

At the CLI/MCP boundary, a proposed request includes the versioned $Q$, evidence references, permitted worker operations, resource limits and logical attempt identity. A result retains the original request alongside the graph/assumption decisions, estimator or refusal, scientific diagnostics, code/artifact hashes and complete tool outcomes. Logs should support reconstruction of why a semantic choice changed; a narrative summary cannot replace the underlying bindings. Generated code runs only in an explicitly qualified worker scope, and a successful interface call grants no scientific acceptance. No repository integration is implemented by this design discussion.~\citep{exec25_cow_reference}

Known repair routes include explicit semantic fixtures, independently derived causal oracles, versioned source custody, reproducible numerical tests and complete cost accounting. Unavailable journal originals, full underlying-study inspection and task-specific contamination evidence remain bounded obligations. Unknown bank assignment truth, unrestricted latent identification and universal transfer of agent accuracy cannot be settled by this public survey. Modular tools make the process inspectable; scientific acceptance still depends on the target, evidence and validation law.~\citep{exec25_cow_reference}


% SOURCE docs/release/2026-10-29/source/docs/writing-repair/2026-10-29/src/b1a07f2bb55d6279.tex CHAR 0
\section[Independent causal and scoring oracles]{Evaluating assistance when the causal answer is known}
\label{sec:exec25-bor}

A public benchmark can justify a prototype only if its reference answer represents the scientific decision we intend to test. An agent can reproduce a benchmark label while answering a different causal question; it can also be penalized for a justified refusal when reference metadata has lost its uncertainty. The retained CLadder and CausalDS papers motivate complementary formal-question and data-science tasks. Their published results remain historical reports under their original versions. Here we inspect two previously identified source issues and derive independent finite witnesses. These specify qualification requirements for the proposed product; they do not report execution of either benchmark or invalidate every item.~\citep{v12_la_2312_04350,v12_la_2607_08093,exec25_bor_sources,exec25_bor_reference}

\subsection{An instrument does not by itself answer the population decision}

In the retained CLadder source, the instrumental-variable graph has an unobserved variable $V_1$ affecting treatment $X$ and outcome $Y$, an instrument $V_2$ affecting $X$, and $X$ affecting $Y$. The corresponding average-treatment-effect metadata contains the ratio of the outcome difference across instrument values to the treatment difference. The question generator separately calls the structural model's treatment-effect bounds. These are distinct parts of the source path: the displayed ratio does not establish that every generated item is wrong, and a structural-model calculation does not establish identification from the information supplied to a respondent. Qualification must bind the graph, parameters, visible information, target and answer-producing path for each admitted item.~\citep{exec25_bor_sources}

We can test the graph-level claim without executing the repository. Let $Z=V_2$ be an independent fair binary instrument and let the latent variable $U=V_1$ satisfy $\Pr(U=0)=1/4$. Set treatment to $X=(1-U)Z$. Thus the quarter of units with $U=0$ comply with the instrument and the remaining units are never treated. Treatment is monotone in $Z$, the first stage is nonzero, and the instrument has no direct effect on the outcome. Consider two outcome mechanisms:
\[
Y_A(x,u)=\begin{cases}x,&u=0,\\1-x,&u=1,\end{cases}
\qquad
Y_B(x,u)=\begin{cases}x,&u=0,\\1,&u=1.\end{cases}
\]
These are stipulated binary structural models with the same graph restrictions. Both mechanisms agree on every outcome observed under treatment assignment: never-treated units have $X=0$ and $Y=1$, while compliers have $Y=X=Z$. The entire observed distribution of $(Z,X,Y)$ is identical across the models, not merely one moment. The independent enumerator verifies every joint probability exactly.~\citep{exec25_bor_reference}

In either model, $E[X\mid Z=1]-E[X\mid Z=0]=1/4$ and $E[Y\mid Z=1]-E[Y\mid Z=0]=1/4$. The ratio is consequently 1. The population effects nevertheless differ:
\[
\tau_A=\tfrac14(1)+\tfrac34(-1)=-\tfrac12,
\qquad
\tau_B=\tfrac14(1)+\tfrac34(0)=\tfrac14.
\]
The ratio equals the complier effect in these constructions. It does not equal either population effect. Because the observed laws coincide while the population effects have opposite signs, no rule based only on that law can determine the correct population-action sign for both models. Randomization, exclusion, monotonicity and relevance do not remove the missing never-treated counterfactual here. Additional effect restrictions, a different target, or further information must be declared rather than silently supplied by a language model.~\citep{exec25_bor_reference}

For a stipulated banking translation, suppose an outcome unit represents USD100 incremental recovery and the decision covers 10,000 comparable accounts. Ignoring action costs only for this illustration, the population increments are USD$-500{,}000$ and USD250,000. Identical observed instrument data therefore support incompatible deployment decisions under compatible structural models. This is a finite decision witness, not a bank estimate. A prototype can return the ambiguity, offer a complier target, or identify additional counterfactual information needed for the decision.~\citep{exec25_bor_reference}

\subsection{Unknown reference truth is a separate state}

The retained CausalDS identification helper represents an unavailable or failed backend with a nullable identification value. Its population-ATE path also excludes an IV-only estimand from accepted nonparametric methods. These are useful distinctions. The grading helpers do not apply the nullable state uniformly: missing ATE identification metadata defaults through a boolean test to not nonidentifiable, while an explicit null value in that same test is negated to true; the inspected counterfactual helper tests whether the value is exactly false. Thus missing, null and false require separate qualification cases. This concerns the inspected helper branches, not a measured frequency in published results. The complete caller and data-validation path still needs native verification.~\citep{exec25_bor_sources}

We propose an explicit oracle state $I\in\{\mathrm{identified},\mathrm{not\ identified},\mathrm{unknown}\}$. Only literal true and false decisions with their scientific evidence enter the first two states; missing values, nulls, backend failures and malformed metadata remain unknown. Unknown tasks enter a quarantine denominator and do not receive fabricated correctness labels. For a known nonidentified target, an intentional refusal is correct and an unsupported estimate is not. A malformed or missing response remains invalid rather than becoming an intentional refusal. For an identified target, a valid estimate proceeds to numeric scoring and a refusal is recorded separately. This prevents infrastructure failure from becoming scientific impossibility or success.~\citep{exec25_bor_reference}

The source already distinguishes intentional null-equivalent responses from invalid answers in its abstention helper; our contract preserves that useful separation. It adds a distinct treatment of uncertain reference truth. A CLI or MCP response should carry target, identification state, reason, source evidence, backend status and response validity separately. Review may resolve an unknown state, but a wrapper must not resolve it by replacing null with a boolean.~\citep{exec25_bor_sources,exec25_bor_reference}

\subsection{Finite intervals and transparent losses}

The inspected ATE interval branch converts endpoints to floating-point values and rejects reversed endpoints. Its local ordering guard alone does not reject every nonfinite value. A comparison of NaN with a finite upper endpoint does not report the lower endpoint as greater, while subtracting NaN produces a nonfinite width. Infinite endpoints can cover a finite truth with infinite width. These arithmetic facts do not establish that every full entry path accepts such values; they establish why public qualification must test parser, validators and scorer together.~\citep{exec25_bor_sources,exec25_bor_reference}

For a finite truth $\theta$, endpoints $l\leq u$, and $0<\alpha<1$, our independent reference uses
\[
S_\alpha(l,u;\theta)=(u-l)+\frac{2}{\alpha}
\{(l-\theta)_+ +(\theta-u)_+\},
\]
where $(a)_+=\max(a,0)$. Coverage, width and loss are distinct outputs. At $\alpha=1/20$ and $\theta=0$, $[-1,1]$ has loss 2 and covers the truth; $[1,2]$ has loss 41 and misses it. The authored validator rejects nonfinite endpoints or truth, invalid confidence levels, reversed endpoints, boolean values masquerading as numbers, and nonfinite arithmetic after finite inputs. Two finite large endpoints can overflow on subtraction. These fixtures are a proposed scientific contract, not a calibration theorem or native conformance result.~\citep{exec25_bor_reference}

The retained CausalDS aggregator reports both uncapped and capped normalized interval-loss means, with separate coverage and width summaries. We retain that distinction rather than claiming the uncapped statistic is absent. Consider 100 stipulated losses for A: 99 zeros and one loss of 1000. B has 100 losses of 1. The raw means are 10 and 1; capping each loss at 10 changes them to 0.1 and 1 and reverses the ranking. A capped mean is a different objective. Neither a low capped score nor high coverage alone establishes useful bank precision or tail-risk control.~\citep{exec25_bor_sources,exec25_bor_reference}

\subsection{A bounded public qualification program}

The proposal can specify a concrete trial without expensive bank evidence. Freeze source, dependency and task identities; construct independent causal oracles; retain separate identified, nonidentified and oracle-unknown partitions. Include observationally identical IV worlds, zero first stages, missing and nullable metadata, invalid responses, finite and nonfinite intervals, arithmetic overflow and capped-loss ranking witnesses. Hold out structures and mechanisms appropriate to the intended generalization claim. Record every submitted request, completed attempt, refusal, invalid response, quarantined oracle and cost. Report conditional scientific accuracy together with those denominators so exclusion of unknown tasks cannot silently inflate apparent coverage.~\citep{exec25_bor_reference}

Adoption remains conditional on matching the full native generator and scoring pipeline to these fixtures, resolving per-item visible information, and reviewing exact artifact rights and dependencies. The mathematical witnesses resolve whether these particular graph and arithmetic pitfalls are possible. They do not close full benchmark qualification. Nor do they supply a universal repair for unrestricted hidden-state identification: where admissible models imply different decisions from identical observations, the proposed product must obtain information, narrow assumptions or return the unresolved set.~\citep{exec25_bor_reference}


% SOURCE docs/release/2026-10-29/source/docs/writing-repair/2026-10-29/src/22103b0d0b8a5b96.tex CHAR 0

\section{Sensitivity and uncertainty are diagnostics, not a causal certificate}
\label{sec:exec23-sua-review}
An agent translating a banking request should notice when harmless rephrasing
changes its proposed variables, graph or estimand. Stability nevertheless does
not establish that a graph is correct. A consistently wrong interpretation can
survive every sampled paraphrase. Conversely, a genuinely different instruction
should sometimes change the output. A diagnostic must specify which
perturbations preserve the business decision and what error it detects.

Hiremath and Hiremath propose sensitivity--uncertainty alignment (SUA) as a
combined diagnostic and training objective. Their version-one preprint defines
an expected divergence between output distributions under perturbation,
subtracts a multiple of predictive entropy, and proposes abstention when the
score is high. It also claims risk and calibration guarantees and reports
experimental tables. Its diagnostic idea does not make those guarantees or
comparisons established premises for this proposal. \citep{exec23_sua_original}

\subsection{What the score measures}
Write $q_x(y)$ for the model distribution, $p_x(y)$ for the true label law,
and $\Pi(\cdot\mid x)$ for a declared perturbation distribution. The
operational quantities are
\[
 S(x)=\mathbb E_{x'\sim\Pi(\cdot\mid x)}D(q_x\Vert q_{x'}),
 \quad U(x)=H(q_x),\quad A_\lambda(x)=S(x)-\lambda U(x).
\]
The supremum $S_\infty(x)=\sup_{x'\in\Delta(x)}D(q_x\Vert q_{x'})$
is different. Generally $S(x)\leq S_\infty(x)$; full support does not reverse
this direction. The source acknowledges it in Remark 3.3 but later substitutes
the expectation into an upper bound involving the supremum. Its entropy-gap
quantity $|H(q_x)-H(p_x)|$ also differs from probability calibration and
accuracy. \citep{exec23_sua_original}

Let true binary label probabilities be $(9/10,1/10)$ and model probabilities
$(1/10,9/10)$. Entropies agree, but the model's modal prediction has true
error $9/10$ and total-variation distance from truth is $4/5$. If its
probabilities remain unchanged across paraphrases, sensitivity is zero too.
This is an authored mathematical example, not a measured LLM failure.
\citep{exec23_sua_counterexamples}

Latent interpretation needs equal care. Two equally likely interpretations
can deterministically map to the same output label. Interpretation entropy
is then $\log2$ while true label entropy is zero. The source's slack
$\eta=|H(Y\mid x)-H(Z\mid x)|$ is $\log2$ and can remove any positive
conclusion. Deterministic mapping alone does not imply zero slack: the output
must also distinguish interpretations. The source's triangle-inequality bound
can remain formally valid but uninformative.
\citep{exec23_sua_original,exec23_sua_counterexamples}

\subsection{Counterexamples to the stated risk guarantees}
\label{eq:exec23-sua-counterexample}
The source initially evaluates risk under the model's own distribution at its
modal prediction. That differs from true banking error. Even at this level,
the claimed entropy subtraction fails. Take binary labels, zero--one loss,
total variation and a constant fair model equal to the true law. A singleton
latent interpretation makes the model exactly equal its interpretation law;
the collapse gap is zero. Choose the allowed concave nondecreasing function
$\psi(h)=0$, $L_D=1$ and $\lambda=1$. The assumptions and condition
$\lambda H\geq\psi(H)$ hold, but equation (11) would imply
\[
 \frac12=R^{\mathrm{rob}}(x)
 \leq R(x)+L_D A_\lambda(x)+\kappa(x)
 =\frac12-\log2<0.
\]
This refutes the stated general bound without a rare adversarial input.
\citep{exec23_sua_original,exec23_sua_counterexamples}

Equation (10) also fails in its stated class of bounded losses. Set
$\ell(y,a)=1/2$ for every label and action, retain the fair model and true
law, and use $\psi(h)=(1-\exp(-h))/2$. This function is nondecreasing,
concave, zero at zero and never exceeds the constant expected loss. The
global entropy--risk assumption therefore holds. The total-variation
Lipschitz assumption holds with $L_D=1$, although its minimal constant is
zero. Sensitivity and collapse gap are zero, $H=\log2$, and $\lambda=1$
satisfies the condition. The claimed upper bound is $1/4$, while actual
self-risk is $1/2$. This is a counterexample to the general bounded-loss
claim, not a proposed banking loss. \citep{exec23_sua_counterexamples}

The proof text reveals the algebraic difficulties. Adding and subtracting
$\psi(H)$ does not permit dropping the positive addend from an upper bound.
The condition $\lambda H\geq\psi(H)$ gives
$-\lambda H\leq-\psi(H)$, opposite to the direction used to justify a
larger right-hand side. Replacing a supremum by a smaller expectation also
tightens an upper bound without justification. Its first step compares the
original distribution's loss at a different modal action with loss at its own
modal action in the wrong direction for zero--one loss. A valid narrow
statement is possible, but these steps do not prove the entropy improvement.
\citep{exec23_sua_original}

A second finite example isolates expected versus worst perturbations. At
the reference input the model is a point mass; at one adverse perturbation
it is fair. Give that perturbation probability $1/100$ and the reference
$99/100$, so the two-point law has full support. Expected total variation is
$1/200$, whereas its supremum and worst self-risk are $1/2$. Small average
sensitivity need not certify worst-case behavior.
\citep{exec23_sua_counterexamples}

Boundedness alone also does not supply the asserted linear constant for any
$f$-divergence. For a binary comparison near equality, a fixed-action
zero--one loss difference is $\epsilon$, but
\[
 D_{\mathrm{KL}}((1/2,1/2)\Vert(1/2+\epsilon,1/2-\epsilon))
 =-\tfrac12\log(1-4\epsilon^2)
\]
is of order $\epsilon^2$. Their ratio diverges. The chosen divergence and
its actual modulus belong in the assumptions.
\citep{exec23_sua_counterexamples}

\subsection{A narrow valid statement and the selection denominator}
For finite labels and zero--one self-risk define $r(q)=1-\max_y q(y)$.
Each singleton probability differs by at most total variation; hence
\[
 |r(q)-r(q')|\leq D_{\mathrm{TV}}(q,q').
\]
Worst self-risk is consequently bounded by $r(q_x)+S_\infty(x)$ when
sensitivity uses total variation. Average perturbed self-risk is bounded by
$r(q_x)+S(x)$. These are different statements, with no unproved negative
entropy term. Neither equates self-risk with error under $p_x$, and neither
identifies an intervention effect. This short derivation is project work.
\citep{exec23_sua_counterexamples}

Selection introduces another denominator. Proposition 4.8 and Appendix A.5
replace risk conditional on the accepted set with unconditional risk,
requiring that accepted inputs be no harder on average. A score threshold
does not establish that condition. In an exact example, one tenth of inputs
have a stable fair model, score $-\log2$ and self-risk $1/2$. The remainder
have a deterministic reference model, positive perturbation sensitivity and
zero baseline self-risk. Threshold zero accepts only the first group.
Coverage is $1/10$, conditional self-risk $1/2$, and unconditional baseline
self-risk $1/20$. The conditional term cannot be replaced by the unconditional
one from the score rule alone. \citep{exec23_sua_original,exec23_sua_counterexamples}

A prospective banking study should freeze task, model, perturbation generator,
threshold and abstention policy, then evaluate accepted answers against
independently adjudicated labels. Threshold selection needs separate data or
an explicitly justified simultaneous procedure. Count accepted and rejected
cases, adjudication failures, semantic changes and clarification costs
separately. Dependent labels or labels supplied by the same model require a
corresponding evaluation law, not an unexamined independent-sample interval.

\subsection{What can be retained from the empirical proposal}
The paper reports accuracy, robustness, calibration and selective accuracy for
broadly described QA, NLI and classification tasks. These are source-reported
numbers. The inspected text does not bind them to exact model checkpoints,
named dataset releases and splits, sample counts, seeds or run artifacts.
Main and appendix comparisons also need consistent baseline and subset
definitions before adoption. This review neither treats tables as reproduced
measurements nor infers fabrication. Missing empirical identities remain
explicit acquisition obligations. \citep{exec23_sua_original}

The source acknowledges perturbation dependence, unobserved latent ambiguity
and a classification focus. Its taxonomy allows meaning-changing inputs with
unchanged outputs to have no SUA signal. For causal work, changing the target
from borrowers to applicants, or default before prepayment to default if
prepayment were prevented, can change the estimand. A semantic reviewer must
preserve these distinctions before output stability is meaningful.
\citep{exec23_sua_original}

Retain sensitivity and entropy as separate logged diagnostics; a combined
score remains a candidate for held-out comparison. Do not use the inspected
generic bounds to certify an intervention diagram, causal estimate or truth.
A matched study can test error detection at fixed coverage and total cost
after fixing perturbation semantics and empirical identities. That research
route is executable; the general certification is rejected as stated.


% SOURCE docs/release/2026-10-29/source/docs/writing-repair/2026-10-29/src/45dcf0b573ca8efb.tex CHAR 0
\section{Evaluating the evaluator: compositional bias and causal acceptance}
\label{sec:exec24-judge-bias}

A bank using an agent to translate business questions into causal studies also
needs a way to review the agent's work. An evaluator may compare a proposed
diagram with a narrative, check whether a report uses the supplied evidence,
or judge an explanation of uncertainty. These are useful tasks, but the
evaluator can itself ignore evidence, reward fluent unsupported answers or
react to irrelevant presentation changes. The relevant question is whether
the evaluation changes for the right reason and remains anchored to an
independently specified scientific standard.~\citep{exec24_judgebias_reference}

\citet{exec24_judgebias_original} study compositional bias in multimodal
language models used as judges. The inspected original is arXiv
2604.18164v3, dated 23 April 2026, before the research cutoff. Its title and
author order match the inherited conference-work metadata, but the final
conference PDF is not asserted to be byte-identical. The earlier v1 original
is retained separately; selected methods in the later revision, rather than
an unlabelled mixture of versions, support the discussion here.
The study concerns pointwise vision-language evaluation. Its empirical
results are not measurements of banking causal identification, diagram
validity, or the proposed library's performance.
\citep{exec24_judgebias_original}

\subsection{Two desired responses to two different changes}

The paper represents a judging input as a query $Q$, an image $I$ and a
candidate response $R$, with a score $y=f_\theta(Q,I,R)$ on a scale from one
to $S$, taking $S=10$ in its experiments. A perturbation produces
$(Q',I',R')$ and a second score $\widehat y$. The distinction between
perturbations is substantive. Removing necessary evidence or mismatching a
response with its evidence should reduce an otherwise favorable score;
changing a semantically irrelevant aspect should preserve it.
\citep{exec24_judgebias_original}

For the first type, the paper defines Bias-Deviation as
\[
 \mathrm{BD}=E\left[\frac{(y-\widehat y)_+}{y-1}\;\middle|\;y>1\right],
 \qquad (a)_+=\max(a,0).
\]
The condition excludes an original score of one because there is no room
for a further decrease. For the second type, Bias-Conformity is
\[
 \mathrm{BC}=E\left[1-
 \frac{|y-\widehat y|}{\max(y-1,S-y)}\right].
\]
The denominator normalizes by the largest available movement from the
original score. These formulas apply to different perturbation families;
using a high score from one as an answer to the other changes the target.
The study groups nine bias types into integrality, congruity and robustness.
Its method therefore requires a defensible prior judgment about whether
each perturbation changes the relevant meaning.
\citep{exec24_judgebias_original}

This requirement is especially consequential for causal diagrams. Moving a
node on the page while preserving every label and endpoint mark can be
irrelevant to the graph's meaning. Removing an arrowhead, deleting a variable
that blocks a back-door path, or moving a measurement across the decision
time is not. A generic visual transformation that is harmless in one task
can obscure an arrowhead or minus sign in another. The proposed bank
evaluation must validate the transformation against the graph or numerical
object before classifying it as meaning-preserving.
\citep{exec24_judgebias_reference}

\subsection{What the empirical design establishes}

The selected paper reports 1,804 evaluation pairs derived from 29 source
benchmarks, with four core task categories and twelve domains. It synthesizes
queries, filters them and uses human review to enforce requirements including
the need to consult the image and consistency with the specified task.
The appendix describes twelve annotators and their selection criteria. This
construction is more informative than a claim that random visual changes
automatically preserve or destroy meaning, but the resulting population is
still the curated benchmark population.~\citep{exec24_judgebias_original}

The study describes 26 judge models, including proprietary, open and
specialized critic families, and evaluates reasoning modes where available.
The main-result caption reports three experimental runs and supporting
variation across runs and samples. The stated token budget is 16,384, with
specified reasoning settings and other parameters left at model defaults.
Those operational choices are part of the inherited result. Repeated runs
help reveal variation for these cases; they do not manufacture independent
new business questions. A new comparison should preserve task identities
and report both within-case variation and variation across distinct cases.
\citep{exec24_judgebias_original,exec24_judgebias_reference}

The paper reports differing strengths across the perturbation families and
does not reduce reliability to model size or a single reasoning switch.
Its prompt interventions likewise produce different trade-offs across
models and bias types. The appropriate proposal use is to adopt the
distinction between relevant sensitivity and irrelevant instability as a
test-design principle, and to reference the observed multimodal failure
modes. Selecting a banking judge solely from its aggregate benchmark rank
would require a transfer argument and new evidence.
\citep{exec24_judgebias_original,exec24_judgebias_reference}

\subsection{Why stability is not correctness}

A constant evaluator illustrates a limitation already acknowledged in the
paper's discussion of supporting sample variance. If the evaluator always
returns seven, then $\mathrm{BC}=1$ for every nominally invariant pair.
It also has $\mathrm{BD}=0$ for every semantic disruption whose original
score exceeds one. Neither statistic by itself establishes accuracy against
a scientific oracle. In an authored four-case set containing two valid and
two invalid causal claims, accepting all four gives perfect score stability
but two false acceptances.~\citep{exec24_judgebias_original,
exec24_judgebias_reference}

Bias-Deviation has its own denominator that must remain visible. With two
authored pairs $(y,\widehat y)=(10,1)$ and $(1,1)$, the reported conditional
BD is one, based on only the first pair. The second pair is not evidence of
a successful response to the perturbation; it is outside the metric's
eligibility set. A bank comparison should publish the total number of pairs,
the number eligible for BD, their task composition and the reasons for
exclusion. It should not convert a conditional metric into an all-request
success probability.~\citep{exec24_judgebias_reference}

Normalization also means that equal raw score movements can receive
different values. A decrease from ten to eight contributes $2/9$ to BD;
a decrease from three to one contributes one. Both are two-point changes.
This is a consequence of the stated normalization, not an arithmetic
defect. It shows why base-score distributions and score-scale interpretation
matter when comparing judges. For causal acceptance, the product additionally
needs an externally defined distinction between a valid identified result,
a justified identified set, a justified refusal and an unsupported claim.
\citep{exec24_judgebias_reference}

\subsection{Abstention requires both conditional quality and coverage}

The paper's abstention-aware variant permits an N/A response, excludes that
response from the numerical metric and reports its proportion separately.
This is a clearly different estimand from a score averaged over every
submitted case. Its methodological implication is to preserve both the
quality of retained judgments and the coverage of the evaluator.
A scientific refusal is not automatically an evaluator failure, but
removing it from every denominator would obscure the service's usefulness.
\citep{exec24_judgebias_original}

Let $J$ indicate that an evaluator returns a scored judgment and $C$ indicate
that the judgment agrees with a prespecified oracle. Then
\[
 P(C\cap J)=P(C\mid J)P(J).
\]
In an authored 100-case example, 50 cases are scored and 45 of those agree
with the oracle. Conditional agreement is 90 percent, while the fraction
of all requests receiving a correct scored judgment is 45 percent. If
some unscored cases represent appropriate refusal, their contribution to
utility should be specified separately rather than labelled either correct
or incorrect by an undocumented convention.
\citep{exec24_judgebias_reference}

An economic loss makes the choice concrete. Suppose the two false
acceptances in the four-case example each expose a stipulated 250,000 USD
loss, while expert review costs 1,000 USD per case. The false-acceptance
loss is 500,000 USD and reviewing all four costs 4,000 USD. These are
hypothetical consequences, not measured bank losses. They demonstrate why
a threshold tuned to an average evaluator score need not minimize decision
loss. With a separate hypothetical expenditure of 15,000 USD for the
100-case exercise, cost is 150 USD per submission, 300 USD per scored case
and $1{,}000/3$ USD per correct scored judgment.
\citep{exec24_judgebias_reference}

\subsection{A bounded route into the proposed product}

The proposed first prototype uses public or generated causal problems with
known structural assumptions and exact refusal oracles. It constructs
paired reports and diagrams, verifies which transformations preserve the
underlying scientific object, and keeps a held-out set of business-to-causal
translations. It then evaluates sensitivity, stability, absolute correctness,
appropriate refusal, coverage and full cost together. The unit of analysis
is the logical case; for example, three attempts on each of 100 cases yield
300 attempts but only 100 distinct cases. A repeated prompt is not a new
independent validation sample.~\citep{exec24_judgebias_reference}

For integration, an evaluator request should carry the causal question,
graph artifact, evidence identities, perturbation class, rubric version,
model configuration and independent oracle where one exists. Its response
should preserve the score, explanation, abstention, failure and cost fields
through the library, CLI and MCP boundary. An LLM explanation may help a
reviewer locate a disagreement, but must not overwrite the engine's
identification or uncertainty status. Scientific owners decide the target
and acceptance loss; protocol owners preserve the evidence and denominator.
\citep{exec24_judgebias_reference}

The inspected study is a useful reference and a candidate source of test
designs. Dataset and code licensing, model access, exact artifacts and
native reproducibility still require separate inspection before direct
adoption. Porting the two explicitly defined metrics is a bounded task,
whereas claiming that a multimodal benchmark qualifies causal reasoning
requires evidence the study does not supply. Its stated limits include
vision-language, pointwise scoring and a focus on multimodal biases rather
than every social or domain-specific bias. The selected preprint supports this comparison; the final version, package execution, independent validation and applicability to banking still require evidence.
\citep{exec24_judgebias_original,exec24_judgebias_reference}


% SOURCE docs/release/2026-10-29/source/docs/writing-repair/2026-10-29/src/b30a0dd0a83384c6.tex CHAR 0
\section{Conflicting modalities require correctness evidence and a separate causal test}
\label{exec23:conflict-scope}

A bank analyst may encounter a chart, its caption and a service narrative that
support different answers. Zhu and colleagues study an adjacent problem in
\emph{Unraveling Cross-Modality Knowledge Conflicts in Large Vision-Language
Models}. The inspected original is arXiv version 1, dated 4 October 2024;
the archive's later version 2 is a separate edition. Their text and image
question-answering experiments motivate checks for inconsistent representations,
but do not establish causal identification or bank decision performance.
\citep{exec23_cross_modality_conflict}

Let $T$ and $V$ be textual and visual answer labels for the same task and let
$G$ be its gold label. Define flip rate $F=P(T\ne V)$ and accuracies
$a_T=P(T=G)$ and $a_V=P(V=G)$. The paper proposes a lower bound on its
knowledge-conflict rate of the form $\mathrm{CR}\geq F-(a_T-a_V)$, attributing
a net accuracy gap to performance-gap samples. That attribution is an
additional premise about unobserved knowledge, rather than an identified
consequence of observing predictions.~\citep{exec23_cross_modality_conflict}

Partition the observations by the correctness of each answer. Cases where
text alone is correct and cases where vision alone is correct both flip.
Both-correct cases do not flip; both-wrong cases may flip. Subtracting the
accuracy difference gives the exact prediction identity
\begin{equation}
 F-(a_T-a_V)=
 2P(T\ne G,V=G)+P(T\ne G,V\ne G,T\ne V).
 \label{exec23:conflict-identity}
\end{equation}
This is a project derivation about answer labels. It supplies no observation
of stored knowledge and no causal explanation of the errors. In a four-item
authored example with gold answers AAAA, text AABB and vision BBAA, each
accuracy is one half, yet every answer flips and the right side equals one.
With text and vision both BBBB, agreement is perfect and both accuracies are
zero. Agreement, correctness and causal validity therefore need distinct
evaluation records.~\citep{exec23_context_reference}

\subsection{The decoding proposal and what the printed results establish}
\label{exec23:conflict-decoding}

The paper reports that confidence, confidence shifts and minimum-variance
strategies did not reliably choose correct answers in the tested conflicts.
Its dynamic contrastive decoding (DCD) uses answer confidence to scale
log-probabilities and subtracts the less confident modality from the more
confident modality. It requires access to the relevant probabilities and
both modality paths. It is a proposed inference-time answer rule, not a
causal intervention on the bank's data-generating process. In particular,
the earlier finding about confidence limits qualifies the heuristic used
to select the subtraction direction.~\citep{exec23_cross_modality_conflict}

Table 4 reports LLaVA-34b accuracy rising from 80.99 to 83.35 on ViQuAE and
from 66.02 to 68.14 on InfoSeek: gains of 2.36 and 2.12 percentage points,
averaging 2.24 points. These are author-reported cells, not reproduced model
runs. The same table contains an InstructBlip-7b InfoSeek tie at 50.53, so
the printed claim of strict improvement across every case is too broad.
There is also a baseline inconsistency: Table 2 prints 82.46 for LLaVA-34b
ViQuAE textual accuracy, whereas Table 4 prints 80.99. Neither value should
silently replace the other.~\citep{exec23_cross_modality_conflict,exec23_context_reference}

Selected Table 5 cells require a stronger arithmetic qualification. Against
the displayed visual baseline accuracy 53.26, the reminder and answer-conflict
values 53.99 and 54.58 differ by $+0.73$ and $+1.32$ percentage points;
the table instead prints $-1.66$ and $-1.07$. Against the displayed recognized
accuracy 58.11, values 57.25 and 58.51 differ by $-0.86$ and $+0.40$;
the printed changes are $-2.53$ and $-1.27$. These four inconsistencies do
not resolve which underlying experiment or baseline is intended. The selected
InfoSeek accuracy changes from 22.11 to 21.25 and 20.23 do match the printed
$-0.86$ and $-1.88$. Preserve the raw cells and their edition while seeking
data or clarification; do not use the inconsistent signs to rank a mitigation.
\citep{exec23_cross_modality_conflict,exec23_context_reference}

\subsection{A bounded route from this literature to a banking prototype}
\label{exec23:conflict-prototype}

The paper's public repository is a candidate for inspection. Its code, data,
license, environment and raw results have not been acquired or qualified here.
The proposal can nevertheless specify useful tests using public or authored
materials. Construct matched chart, table and narrative versions of the same
known-law problem, deliberately introduce a unit change, stale date, swapped
axis or unsupported causal caption, and retain a source-backed gold answer
for each representation. Evaluate answer correctness and conflict detection
separately; count abstentions and incorrect agreement. Require a causal
question to pass the separate target, graph and identification checks developed
above even when all representations agree.

Compare a single-modality baseline, source inspection, explicit disagreement
reporting and, where logits and rights permit, a reproduced DCD rule. Charge
for both modality passes, repeated confidence samples and human adjudication.
The inspected appendix's ten prompt or dropout samples concern its confidence
analysis; they are not a general guarantee about uncertainty or an automatic
ten-fold price for DCD. Freeze tasks, corruptions, costs and denominators before
evaluation. A generated monetary target can expose a wrong decision under a
declared law, while authentic bank utility remains a later experiment. This
sequence makes the expensive study conditional on a demonstrated prototype merit.


% SOURCE docs/release/2026-10-29/source/docs/writing-repair/2026-10-29/src/0288d67fd37fd73b.tex CHAR 0
\section{Shared evidence and scientific validation in a causal assistant}
\label{exec23:admit-scope}
A bank analyst may ask whether a fee waiver prevents delinquency, while several
agents retrieve policy manuals, propose confounders and restate the same
customer-service memorandum. Repetition can make one unsupported assumption
appear to have several independent witnesses. A second failure occurs when
a fluent derivation is accepted because a validator assigns it a low risk
score. Neither successful storage nor a reassuring score establishes that the
waiver is identifiable from the available data. The useful product therefore
needs distinct contracts for admitting evidence, proposing a graph, validating
a mathematical argument and releasing a decision. Two recent papers expose
both useful mechanisms and the assumptions that such contracts must retain.

\subsection{What an admission benchmark actually measures}
\label{exec23:admit-memory}
\citet{exec23_epistemic_admission} distinguish a static admission benchmark
from a running shared-memory environment. The static construction imports
public verdicts, adds authored agent and consumer settings, and scores a fixed
gold action. Its 1,200 episodes are divided into 600 training, 200 validation
and 400 test episodes. Agreement with the admission specification is the
target; it is not an empirical measure of emergent propagation. The live
construction uses six agents for six rounds and 100 configurations per model
family, with additional diagnostic and truth-by-copy sets reported separately.
Every live scenario is authored fiction. This makes the input lineage and
scenario truth available to the evaluator, but does not estimate their
frequency in bank operations.

Only promotion writes to shared memory. Later candidates depend on what
earlier policies admitted, so a false-admission fraction is conditional on a
policy-specific stream. If $N_\pi^F$ false candidates arrive under policy
$\pi$ and $A_\pi^F$ are admitted, the ratio
\begin{equation}
 f_\pi=A_\pi^F/N_\pi^F
 \label{exec23:admit-stream}
\end{equation}
is not a comparison on common candidate inputs. For example, ten admissions
from 100 false candidates give $0.10$, whereas five from ten give $0.50$. The
first policy has the smaller conditional rate and the larger number of
admitted false claims. These are stipulated counts, not the paper's results.
A prototype should therefore use both frozen-candidate replay, which measures
a gate on matched inputs, and a closed-loop experiment, which measures the
whole policy-induced stream. It must not combine their denominators.

The paper's propagation damage counts consumer reads while a false belief is
retrievable. This is simulated exposure, not monetary damage. Assertion
adoption is a separate consumer-output measurement. The reported
$0.97$--$0.99$ adoption of an uncontested false belief is conditional on a
memory-only consumer and a selected pool of 553 probes. It is informative
about that constrained interface, not a bank-wide error probability.
Source-type gating performs well in several authored scenarios but admits
false restatements when the source type supplies the authority signal it
trusts. The lineage oracle rejects copied support yet can admit independently
restated falsehoods. Provenance and truth are different properties
\citep{exec23_epistemic_admission}.

An elementary model makes the distinction precise. Suppose each of three
independent witnesses is correct with probability $p=4/5$. A majority is
correct with probability $3p^2(1-p)+p^3=112/125$. If all three repeat one
root witness, the same majority is correct with probability $4/5$. Copy
detection prevents an unjustified increase in confidence; it cannot improve
the root witness's truth probability by itself. These independent rational
calculations are project examples, not a model fitted to the benchmark
\citep{exec23_admission_conformal_reference}.

\subsection{Truth, lineage and usefulness require separate tests}
\label{exec23:admit-validation}
Input-feed lineage is fixed by the live scenario. For generated claims,
retrieval is logged but the derivation link remains agent-declared and can
be incomplete. A citation may mention a false proposition without endorsing
it, while an uncited sentence may assert the same falsehood. The benchmark
therefore separates assertion grading from provenance tracing. Its assertion
grader receives gold propositions; a deployable admission policy does not.
The distinction must survive any reuse of its evaluation design.

The paper also separates blind human labels from labels revised after seeing
the grader. On the assertion sample, the two blind annotators agree more
closely with one another than either does with the model grader. The consumer
matcher misses rephrased values, and its validation sample is balanced by
matcher category rather than drawn to estimate deployment prevalence.
Reported false assertions and adoption therefore depend on the measurement
procedure. The release should retain raw sentences, blind adjudication,
normalization rules and alternative labels instead of treating one grader as
truth. No evaluated admission policy retracts beliefs although the store
offers retraction actions; correction remains untested, not technologically
impossible. The Mem0 and A-MemGuard comparisons concern particular components
and configurations. They do not license an unrestricted conclusion about
every version or deployment of either system
\citep{exec23_epistemic_admission}.

For the bank prototype, each proposed graph edge should retain the business
assertion, original source, observation date, copying lineage, conflicting
evidence and status of human review. An edge such as ``delinquency causes a
servicing call'' must not become ``a servicing call causes delinquency'' merely
because several agents reproduce the same text. A claim-dependency graph
records support between sentences; an intervention graph records assumptions
about mechanisms. The former can organize evidence for the latter, but cannot
identify an intervention by its own topology.

Usefulness must include abstention. Let $p_C$ and $p_W$ denote probabilities
of a correct and a wrong released answer per evaluated request, let an
abstention have zero stipulated gross value, and let $c$ be mean cost per
request. Under assumed benefit $B$ and loss $L$, the expected net is
\begin{equation}
 V=Bp_C-Lp_W-c.
 \label{exec23:admit-value}
\end{equation}
Take $B=10$, $L=100$ and $c=1$ dollars. A policy with
$(p_C,p_W)=(0.60,0.05)$ has net zero dollars; one with
$(0.30,0.01)$ has net one dollar despite answering correctly less often.
These inputs illustrate a decision rule and are not observed bank values.
Compare admission methods at declared coverage, error costs and workload,
then vary those assumptions. Agreement rates or token counts alone cannot
select the better bank policy~\citep{exec23_admission_conformal_reference}.

\subsection{What graph-structured conformal validation can promise}
\label{exec23:admit-conformal}
\citet{exec23_scientific_conformal} propose progressive scientific validation
with branching and an approximate deducibility graph. Their factuality
definition requires both correct claims and justification from prior context
and scientific ground truth. It does not require a relevant answer. Their
approximate graph has semantic ancestor-connectivity and consistency
properties. Acyclicity, reproducible template matching or a plausible
linguistic causal connective does not alone establish these properties.

There is a sound conditional idea in the process-level argument. Let
$Z_1,\ldots,Z_{n+1}$ be exchangeable complete examples and let $g$ be the same
fixed measurable construction for every example. Then
$(Z_i,g(Z_i))_{i=1}^{n+1}$ is exchangeable because applying the same coordinate
map commutes with permutation. This argument permits rich dependence within
one example. It neither assumes that the claims inside a response are
independent nor manufactures exchangeability across examples. Training a
validator, choosing a threshold on the test cases, changing the generation
policy after calibration, or selecting convenient branches requires a
separate argument. Randomized construction is possible if its randomness is
incorporated with the appropriate symmetric law; empirical repeatability
below one does not verify the paper's deterministic special case.

For a conventional score-coverage construction, a finite calibration sample
uses rank $k=\lceil(n+1)(1-\alpha)\rceil$. If $k=n+1$, an infinity sentinel
may be needed; silently choosing the largest finite calibration score changes
the procedure. With $n=5$ and $\alpha=0.1$, the rank is six. Most importantly,
coverage of a score or prediction set must be linked to the desired event
that the released reasoning is valid. That link requires trustworthy
calibration labels, the correct threshold orientation, a nested filtering
family and the stated semantic conditions. The paper's coherent-factuality
proof invokes those properties in addition to exchangeability. A fallible
validator accepting a false statement with score zero is enough to show why
``accepted prefix'' cannot be substituted for ``true prefix'' in a cascade
prevention argument~\citep{exec23_scientific_conformal,exec23_admission_conformal_reference}.

A marginal guarantee is also different from simultaneous reliability. Five
independent answers, each valid with probability $0.9$, have probability
$1-0.9^5=0.40951$ of at least one invalid answer. Adaptive repeated queries
need a specified selection and error-control contract. This calculation
does not refute a marginal conformal result; it identifies the additional
product question~\citep{exec23_admission_conformal_reference}.

\subsection{Ambiguities to repair before certification}
\label{exec23:admit-repairs}
Several passages prevent direct adoption of the paper as a certificate. Its
main pipeline accepts low risk scores, while its branching prose selects
the highest score above threshold. Its minimum-domain-quantile rule must be
interpreted against an explicit event. For ordinary coverage
$\Pr(S\leq q)$, decreasing $q$ is not conservative: two deterministic
domains with scores and valid own-domain quantiles zero and one each have
own-domain coverage one, but using the minimum quantile zero gives coverage
one half under an equal mixture. In contrast, removing more claims can
preserve an all-retained-claims-valid property when the candidate family is
properly nested and validity hereditary. The counterexample concerns the
ordinary score event, not that distinct filtering theorem. The implementation
must state which event it guarantees and prove the relevant nesting rather
than transport the word ``conservative'' between them
\citep{exec23_scientific_conformal,exec23_admission_conformal_reference}.

The displayed progressive-generation pseudocode increments its attempt
counter only when it is below the maximum. If the last attempt again fails,
the loop condition remains true and no exit or counter increment is shown.
An independent bounded trace with maximum two gives attempt states
$1,2,2,2,2$. This is a finding about the printed algorithm, not evidence that
an uninspected implementation hangs. A prototype should use a bounded
attempt loop, distinguish an incomplete retained prefix from a completed
answer, and return an explicit abstention when its budget is exhausted.
The paper's cycle-to-linear-graph fallback likewise needs semantic testing;
turning a graph into a DAG does not establish the claimed deduction relation.

The reported empirical ranking also needs reconciliation. Table 1 gives
MATH accuracy $48.9$ for the descendant-weighted method and $71.2$ for one
listed comparator, despite its universal-superiority caption. Table 2 changes
the full-system MATH and ScienceQA accuracies, and a partial configuration
reports PhyX accuracy $53.1$ and validity $92.1$, compared with $50.1$ and
$91.7$ for the full system. Main-text and appendix model roles, overall
validity, and some domain figures also differ. These observations justify
requesting exact run manifests, denominators, paired outcomes and frozen
calibration sets before ranking methods. They do not establish misconduct,
nor do they erase the useful design ideas. The local review reproduces
arithmetic and control-flow witnesses only; it does not reproduce the
reported model experiments~\citep{exec23_scientific_conformal,exec23_admission_conformal_reference}.

\subsection{A bounded public-evidence prototype route}
\label{exec23:admit-route}
The proposal can proceed without costly bank experiments. Construct public
or authored business scenarios with declared truth and copying lineage; keep
scoring-only labels out of agent inputs. Freeze separate tasks for translating
a business request, proposing competing intervention graphs, identifying a
query, checking an argument and deciding whether to answer. Replay admission
gates on common candidates, then test the entire interacting pipeline with
independent scenario units. Include paraphrase, conflicting authoritative
text, missing lineage, incorrect graph orientation, unsupported deduction,
cyclic dependency output and exhausted validation budgets.

Fit scoring components on training examples; choose the complete generation,
branching and validation procedure before calibration; evaluate on untouched
examples under the declared sampling law. Record refusals, incomplete answers,
truthful corrections, coverage, resource use and conditional errors separately.
Inspect authentic software, licenses and run artifacts before choosing direct
adoption or a port. The memory benchmark is presently a design reference and
the scientific-validation paper a hypothesis and cautionary comparison.
Neither supplies a ready certificate for causal discovery, bank identification
or universal scientific truth. This staged use of public evidence can establish
whether a prototype merits a later bank study while preserving the scientific
questions that the eventual study must answer.


% SOURCE docs/release/2026-10-29/source/docs/writing-repair/2026-10-29/src/06f70a1512e53b8d.tex CHAR 0
\section{Original factuality guarantees and an agent-facing econometrics comparator}
\label{exec23:proof-scope}
The preceding discussion leaves two different uncertainties: which validity
claim the original conformal arguments support, and whether a unified
agent-facing estimator interface already exists in a form worth reusing.
The original factuality papers and the StatsPAI software paper permit more
precise answers. They narrow the design obligations without supplying an
accepted causal product.

\subsection{A lower guarantee is not a tight calibration interval}
\label{exec23:proof-lower}
\citet{exec23_mohri_factuality} formulate factuality as entailment by a chosen
reference. Their guarantee is relative to the available binary entailment
oracle, which may be a human or an automated judge; it is not a theorem that
the oracle represents external truth. A back-off mechanism progressively
weakens a response and must admit an empty response. Its safe-threshold
score checks every subsequent threshold, not merely one currently accepted
claim. This distinction permits their lower bound without nested entailment
sets. In their Theorem 4.1, nesting is an additional condition for the tight
upper bound. The theorem also restricts the finite-sample error level to
$\alpha\geq 1/(n+1)$ and specifies distinct scores. The infinity-sentinel
example in the preceding discussion describes a conventional extension
outside that finite order-statistic range; it is not a claim that this
published theorem covers $n=5,\alpha=0.1$ as written.

Their practical merger assumption states that a merged output is entailed
if and only if each retained claim is entailed. This is a semantic property
of decomposition and recombination, not stochastic independence of claims.
The exchangeability requirement concerns complete input/reference examples.
Correlated claims inside one example do not by themselves violate that
sampling assumption. The distinction is particularly important when a bank
answer combines accounting identities, a causal assumption and an estimate:
all three may depend on the same underlying records while requiring different
truth checks. Their empirical study uses the first 50 inputs from each of
three datasets, author-provided labels and repeated random half-splits. A
thousand re-splits do not create a thousand independent newly collected
samples. Their limitations explicitly distinguish marginal validity from
per-user, per-domain and fixed-calibration-set guarantees, and warn about
changing input distributions~\citep{exec23_mohri_factuality}.

\citet{exec23_coherent_factuality} make a related but distinct separation.
Their lower bound calibrates the maximum tolerable risk using actual
coherent-factuality labels for candidate subsets. That lower-bound argument
does not require the proxy graph to satisfy every semantic property of an
approximate deducibility graph. Their tight upper bound additionally uses
the approximate-graph properties and a finite first-bad-subgraph threshold.
This qualifies the preceding statement about graph and nesting requirements:
those requirements cannot all be treated as necessary for every lower-bound
construction. They matter for the particular tight-bound proof and for the
shortcut that replaces full subset labels with graph-assisted individual
labels. Genuine subset-level truth information remains the crucial resource.

The paper labels graph-assisted individual assessments as silver and direct
subset assessments as gold. Silver labeling asks whether a claim would be
true if its ancestors were true. Turning these labels into a subset truth
assessment presumes adequate graph structure. The MATH evaluation contains
both label types, whereas FELM contains only silver labels and exhibits
some miscalibration. The authors explicitly say that their model-generated
graph validity cannot be enforced rigorously, report difficulties with long
outputs, and find nonsensical or cyclic biography graphs. Their re-prompted
completions are evaluated separately: observed error reduction is not an
automatic extension of the original filtered-output guarantee
\citep{exec23_coherent_factuality}.

For a prototype, choose one of two auditable services. A full subset-label
service evaluates the actual retained reasoning, with its annotation cost
and subjective rubric recorded. A cheaper graph-assisted service must test
the graph and labeling shortcut against independent subset judgments and
retain the resulting uncertainty. Neither service may identify truth with
a model's confidence score. Fit scorers separately, freeze the complete
pipeline and query population, retain tie rules and empty-output behavior,
and evaluate answer usefulness alongside validity. The original theory
supports these bounded choices; it does not certify unrestricted banking
questions or an adaptively changing assistant.

\subsection{StatsPAI changes the integration comparison}
\label{exec23:proof-statspai}
The published StatsPAI paper describes a unified Python interface spanning
regression, panels, IV, modern difference-in-differences, regression
discontinuity, synthetic control, matching and heterogeneous-effect methods.
Its registry supplies argument/output schemas, limitations, citations,
stability tiers and structured failures for agent-facing discovery, including
possible MCP use. The authors identify the integration layer, rather than a
single superior estimator, as their contribution. Their package depends on
established numerical and econometric libraries and offers optional PyTorch,
JAX and Rust/PyO3 paths. These are source-reported capabilities, not results
of execution in this proposal~\citep{exec23_statspai_paper}.

The comparison is substantive because the proposed library also needs a
common method-discovery and result surface. Before building another broad
wrapper, inspect whether a StatsPAI adapter can preserve the proposal's
estimand, population, intervention, observation clock, uncertainty type,
assumption failures and refusal states. A shared coefficient field is useful
only if it does not disguise different causal targets. The paper limits its
uniform reporting/export promises to mature estimators; auxiliary helpers
have narrower capabilities. Dedicated packages may expose more design-specific
options, and optional compiled kernels add platform and build costs. A
registry count therefore measures interface breadth, not the number of
independently qualified causal methods~\citep{exec23_statspai_paper}.

The paper reports a cross-language parity harness covering more than sixty
estimator modules, with matched input bytes and disclosed tolerances. It also
distinguishes redistributable real teaching data from deterministic simulated
replicas calibrated near published estimates. Those replicas are not exact
replications of the original studies. Its Texas synthetic-control example
reports different local optima and roughly three-percent effect agreement;
a smaller pre-treatment fit error does not by itself establish a more accurate
counterfactual. Its reported 1,000-replication coverage checks concern selected
OLS, difference-in-differences and strong-instrument IV designs. They do not
establish weak-instrument validity or arbitrary bank dependence. Each claim
requires the exact source, data identity, harness, tolerances and output
artifacts before this proposal can report independent reproduction
\citep{exec23_statspai_paper}.

The paper supports conditional adapter/reference use. The following source reviews, beginning in Section~\ref{exec23:statspai-contract-scope}, inspect the cutoff release, schema, error handling and selected runtime helpers. They resolve those specific source questions while leaving native qualification open. A bounded design must still be mapped end to end and compared against an independent reference before choosing a direct dependency, isolated worker or selective port. Reject a route if it loses the
scientific target or silently converts unsupported states into numeric
answers. The paper declares the package MIT licensed and distinguishes it
from a commercial downstream product; package, dependency, dataset and
commercial-use rights still require their own exact records. This route
can avoid unnecessary duplication while keeping native qualification and
bank adoption as later, separately evidenced decisions.


% SOURCE docs/release/2026-10-29/source/docs/writing-repair/2026-10-29/src/69369ed8710421b8.tex CHAR 0
\section{From an agent-facing toolkit to a qualified causal adapter}
\label{exec23:statspai-contract-scope}
The preceding software comparison motivates a more precise design decision.
A banking analyst may ask an agent to estimate the effect of a servicing
intervention, compare alternative estimators, and preserve a result that another
analyst can reproduce. A unified catalogue can reduce the effort of finding and
calling methods. It does not itself determine which customers, intervention,
time horizon or assignment law make the returned number relevant. StatsPAI
provides a concrete integration candidate against which to develop this
separation. The inspected cutoff source declares version 1.39.0 at commit
\texttt{35d5a4fbd05f1b90201c24d806df9b12f966e50d}, dated 7 October 2026.
This is distinct from the paper's archived version 1.23.0; an observation about
one version cannot silently be transferred to the other. The selected eleven
files include the package declaration, result-schema exporter, parity taxonomy,
MCP server, workflow guide and two test specifications.
\citep{exec23_statspai_contract,exec23_statspai_paper}

The package declares Python 3.10 or later, a substantial scientific and reporting
dependency set, optional JAX and PyTorch stacks, and separate command-line and
MCP entrypoints. Its inspected MIT notice includes an attribution condition and
warranty disclaimer. The package declaration also names third-party notices;
those notices, dependency-specific terms and dataset rights require their own
inventory before a chosen distribution. A package-level license observation
is not a conclusion about every optional component or about the separately
identified commercial downstream product. An adapter should pin the source,
selected extras and executable environment together, rather than relying on
lower bounds in a package declaration as a reproducible lock.
\citep{exec23_statspai_contract}

\subsection{Transport shape and scientific admissibility}
\label{exec23:statspai-contract-schema}
There are at least three distinct contracts. A transport contract determines
whether the response can be read. A result contract determines whether the
required fields and diagnostics for a particular method are present. A
scientific contract determines whether the estimand and uncertainty answer the
banking question under stated assumptions. The exported agent-result JSON
schema requires only a string-valued \texttt{method}. It permits method-specific
extensions, and its interval array has no length or order constraint. Its
standard error, probability and observation-count properties do not impose
nonnegativity, the unit interval or integer-valued sample size respectively.
The MCP server's compact advertised output schema is more permissive still:
it accepts any object and refers the client to a fuller result resource.
This supports a heterogeneous catalogue, but acceptance under either schema
must not be promoted into scientific admissibility.
\citep{exec23_statspai_contract}

Independent data-only checks make the boundary concrete. A payload containing
only a method name satisfies the exported schema. A second payload with a
reversed interval $[2,1]$, standard error $-1$, $p$-value $2$, error level $-1$
and observation count $0.5$ also satisfies it, as does an empty interval array.
A numeric method name fails. These are deliberately authored payloads validated
against the downloaded schema; they are not observations that an estimator
emitted such values. The check executes the project's validator and parses
upstream syntax as data; it does not import or execute StatsPAI.
\citep{exec23_statspai_contract_reference}

For a proposed average-treatment-effect adapter, the consumer contract should
require the target population, treatment contrast, units, horizon, information
cutoff and identification assumptions, together with a declared uncertainty
kind. If a finite standard error is applicable, require it to be nonnegative.
If an ordinary confidence interval is applicable, require two ordered endpoints
and an explicit coverage level. A legitimate unbounded identified set must
retain its meaning through a separate representation; turning an infinity into
an apparently missing scalar would erase information. Regression coefficient
tables, causal effects, diagnostic tests and helper responses need distinct
variants. Refusal and partial identification are substantive outcomes, not
malformed estimates to be replaced by a default value. These requirements are
our proposed consumer design, beyond the generic transport schema.

\subsection{Failures, time budgets and reproducible custody}
\label{exec23:statspai-contract-runtime}
The inspected server distinguishes JSON-RPC request failures from failures
inside a tool call. It constructs an execution-error response with an error
flag, structured error fields and a serialized text representation. The
underidentified-IV tests specify retention of diagnostic and alternative-method
information. Those tests are useful contract specifications; their presence
is not evidence that they passed in a local executable environment. Likewise,
a suggested alternative method is a proposal for analyst review, not proof
that it is identified under the original question's assumptions.
\citep{exec23_statspai_contract}

Timeout semantics particularly matter after an interrupted session. The caller
can dispatch an eligible call to a process worker, but a result that cannot be
transferred can fall back to the thread runner. The process-timeout branch
reports a killed worker. The thread branches explicitly report that a data read
or computation may remain active. Eligibility, kill behavior, handle transfer,
queue admission and orphan tracking are delegated to other modules outside this eleven-file inspection. Section~\ref{exec23:statspai-runtime-scope} examines selected additional helpers and narrows these questions; it still does not establish native lifecycle behavior. The supported conclusion is therefore
about the caller's distinct contracts, not a verified guarantee that every
call stops computing when its client stops waiting. An adapter should retain
actual execution mode and worker state, reconcile the original attempt after
reconnection, and allow a retry only after its duplication and resource
consequences are understood. This is the same distinction between an
observation deadline and terminal completion required of commercial workflow
services.
\citep{exec23_statspai_contract}

Output size has a similarly precise boundary. The server applies its budget
to the structured result and then serializes that result again in a text
content block; an image can be appended separately. The stated budget is not
a cap on the entire transmitted message. A structured JSON payload of
262,144 bytes already contributes at least 524,288 bytes when present twice,
before enclosing fields, string escaping or an image. This is an arithmetic
lower bound on the duplicated representations, not a network benchmark.
A causal adapter should measure and limit the complete response and preserve
whether risk details were shortened. The inspected caller records nonfinite
locations, truncation metadata and a capped list of runtime warnings; it calls
separate helpers for the detailed budgeting policy. A small successful message
is not sufficient if omitted diagnostics could change the decision.
\citep{exec23_statspai_contract,exec23_statspai_contract_reference}

The workflow guide recommends saving externally retrieved data and preserving
provider, indicator, query, retrieval date and transformations. It distinguishes
local-file provenance from remote data whose bytes are not hashed unless saved
locally. The caller also distinguishes a standalone replay from a replay that
depends on process-held handles or merely records a call. Thus a replay string
is not a complete reconstruction package. For a banking proposal, the reusable
artifact should contain the selected immutable input, transformation graph,
model configuration, sampling and fold choices, package environment, complete
result and diagnostics. Its permissions and confidential-data custody must be
specified separately. No sibling repository has to be modified to define this
handoff: a consumer can read a versioned artifact and test the declared result
contract in its own environment.
\citep{exec23_statspai_contract}

\subsection{Parity is conditional on the target and the statistic}
\label{exec23:statspai-contract-parity}
The taxonomy explicitly records past problems with inherited alias claims and
circular evidence. It separates estimator entries from loaders, renderers and
other infrastructure, and requires named equivalence evidence before an alias
receives another entry's grade. Its mapping labels externally aligned results
as certified, internally supported results as validated, and other states as
API-stable. These are upstream labels with specific meanings. None is an
independent certification of a bank's causal claim. The source's Python-reference
mapping contains eight entries although the adjacent comment still says two;
the literal mapping, verified without executing the module, prevents the stale
comment from narrowing the apparent cross-package exceptions. Cross-package
comparison should not be described as cross-language evidence when both
implementations are Python.
\citep{exec23_statspai_contract,exec23_statspai_contract_reference}

The inspected parity-boundary tests read committed results. For one RD fixture
they prescribe much tighter point-estimate tolerances than the allowed relative
standard-error difference of seven percent. A DML comparison specifies shared
explicit folds; a synthetic-control comparison distinguishes nonunique weight
solutions and a unique-solution counterpart. These are useful prospective
checks. They do not show that the current installation reproduced R or Stata,
or that a tight agreement on one fixture establishes coverage for another
sampling law. An adapter's evidence unit must therefore name the exact entrypoint,
input bytes, target, statistic, convention, reference version and tolerance.
The guide itself separates stored comparisons from invoking a separately
configured R or Stata runtime.
\citep{exec23_statspai_contract}

A documented non-alias illustrates why matching a method family is insufficient.
The taxonomy distinguishes cohort-size weighting from treated-observation
weighting even when both arise from related difference-in-differences estimators.
Consider our stipulated servicing example with two cohorts of ten customers
each. The first cohort contributes three post-treatment periods and has a
constant effect of USD 4 per customer-period; the second contributes one period
and has effect USD $-4$. The equal-customer average effect at a common evaluation
period is
\[
 \tau_{\mathrm{customer}}=\frac{10(4)+10(-4)}{20}=0,
\]
whereas the available treated-observation average is
\[
 \tau_{\mathrm{observation}}=\frac{10\cdot3\cdot4+10\cdot1\cdot(-4)}{40}=2.
\]
With a stipulated USD 1 per-customer decision cost, the first target gives
net USD $-1$. Subtracting that cost from the second average gives USD 1 but
mixes a different weighting target into the customer decision. Neither
arithmetic calculation repairs the mismatch. No estimator is run here, and
these dollar inputs are authored examples, not measured benefits. The adapter
must bind the aggregation and horizon before deciding whether two software
outputs disagree or simply answer different questions.
\citep{exec23_statspai_contract,exec23_statspai_contract_reference}

\subsection{Adoption decision and a bounded qualification route}
\label{exec23:statspai-contract-decision}
The preferred proposal route is a conditional isolated adapter, with StatsPAI
serving as a source-grounded reference until qualification. Its catalogue,
metadata, structured errors and replay concepts are valuable design assets.
Direct blanket adoption is deferred because the generic schema does not enforce
our scientific contract. The additional helper inspection in Section~\ref{exec23:statspai-runtime-scope} resolves part of the source uncertainty and identifies narrower lifecycle, custody and replay limitations; native guarantees remain unverified.
A selective port may be justified by a specific numerical or deployment need,
but a broad port would add solver, environment and maintenance obligations.
The rejected claim is blanket certification based on registry labels or JSON
validity; rejecting that claim does not discard the underlying package.
\citep{exec23_statspai_contract}

The executable qualification route has five bounded parts. First, inventory
third-party notices, selected dependencies and data rights, and lock one
estimator family and environment. Second, use the following same-revision helper review to specify the process, loader, cache, replay and budgeting tests, and inspect any additional selected path. Third, when a
separately permitted prototype environment exists, reproduce one named estimator
against a qualified reference with identical inputs and folds, preserving full
covariance, diagnostics and refusal behavior. Fourth, exercise stale handles,
missing files, nonfinite values, incomplete risk output, timeout, cancellation
and reconstruction after restart; scientific success must not survive an
incomplete result. Fifth, compare the adapter with a direct mature-package
wrapper on matched public or generated tasks, recording actual operational
cost separately from stipulated banking economics. Failure at one gate gives
a scoped disposition and does not block independent literature or software
review. This develops a testable product proposal while leaving native
implementation, independent acceptance and bank adoption as later evidence.


% SOURCE docs/release/2026-10-29/source/docs/writing-repair/2026-10-29/src/c6dca278318f2e69.tex CHAR 0
\section{Reproducing the calculation and its evidence}
\label{exec23:statspai-runtime-scope}
The conditional adapter decision can now be narrowed by inspecting eight
additional files at the same StatsPAI cutoff revision. These cover third-party
notices, process and thread workers, output budgeting, local data loading,
result and data caches, and replay construction. The additional inspection
resolves specific source questions left open in the preceding discussion.
It does not execute the toolkit or convert source-defined behavior into a
measured reliability claim. The third-party notice reproduces an MIT notice
for the HonestDiD 0.2.8 algorithm port in a named module; the selected dependency
stack, redistribution and datasets still require a complete inventory.
\citep{exec23_statspai_runtime}

\subsection{An observation deadline is not a complete resource contract}
\label{exec23:statspai-runtime-deadline}
Process isolation is opt-in; the default is the thread runner. The eligibility
predicate excludes calls with data handles, selected session tools and pipeline
prefixes. Result handles may be transferred using signed serialized objects,
but a missing or unpicklable result can cause fallback to the thread path.
The HMAC check establishes that a transfer was signed with the shared per-call
key; it is not a scientific validity test or a sandbox for arbitrary trusted
Python objects. The proposed adapter must declare whether a tool is eligible
and record which mode actually executed, rather than infer isolation from a
server-wide option.
\citep{exec23_statspai_runtime}

The process supervisor creates a child, writes its request, then establishes
its deadline. On timeout or cancellation it calls the direct child's kill
method and waits up to ten seconds. After receiving a result it can wait up to
thirty seconds for the child; when stdout closes without an answer, a branch
calls an unbounded wait. These are static control-flow observations, not
observed hangs. The inspected routine does not establish a process-tree
termination guarantee, nor does its clock cover every startup, request-write
or post-response operation. A native qualification should therefore measure
end-to-end elapsed time separately from estimator time and test startup failure,
closed streams, descendants and delayed exit in a permitted environment.
The adapter's own supervisor needs explicit terminal states and resource
ownership, including a response that says termination could not be confirmed.
\citep{exec23_statspai_runtime}

The thread worker signals cancellation and registers unfinished threads as
orphans. Cancellation takes effect at a cooperative checkpoint or when the
operation returns; the response deadline does not kill the thread. In the
inspected configuration reader, the default is 600 seconds and a nonpositive
value disables the timeout. A proposed servicing workflow must budget active
attempts as well as client waiting time. In an entirely stipulated illustration,
a single useful outcome worth USD 10 and an attempt cost of USD 7 gives net
USD 3; two charged attempts delivering that same one outcome give USD $-4$.
This arithmetic does not measure infrastructure prices. It shows why retry
policy belongs in the full workflow comparison and why reconnecting should
first reconcile the original request.
\citep{exec23_statspai_runtime,exec23_statspai_runtime_reference}

\subsection{Resource bounds and immutable inputs}
\label{exec23:statspai-runtime-custody}
The output-budget helper protects headline fields and selected evidence
subtrees, shortens other containers first, and reports incomplete risk details
when risk lists must be shortened. It can report unavoidable overflow when
protected content itself cannot fit. Its scope remains structured content,
so the consumer must budget the whole wire representation separately. This
is a valuable explicit failure state: an oversized protected diagnostic should
lead to an artifact handoff or bounded refusal, not silent deletion of the
scientific reason for rejecting an estimate.
\citep{exec23_statspai_runtime}

The data cache estimates DataFrame memory usage, whereas the general result
cache is count-bounded by default and has optional expiry. Its insertion loop
evicts for byte pressure only while more than one entry remains. Consequently,
the source permits the newest single item to exceed the byte budget. A small
independent model of that guard retains an item of size 12 against a budget of
10. This is not a runtime memory measurement. A strict deployment limit needs
an explicit oversize-item decision before insertion, a clear distinction
between object estimates and process memory, and resource measurements that
include workers and serialization copies. Cache handles also do not supply
durable storage: lineage stops at an evicted parent, and its traversal has a
finite depth. The adapter should retain an external immutable reconstruction
artifact when the scientific record must survive process loss.
\citep{exec23_statspai_runtime,exec23_statspai_runtime_reference}

The local loader resolves paths for an allowlist check, but no roots configured
means unrestricted local paths within the process's own permissions. It then
opens the caller path through format-specific readers. Data-frame caching uses
path, modification time and requested columns; hash caching additionally uses
file size and nanosecond modification time. These are efficient cache keys,
not immutable identities. Two distinct equal-length byte strings have different
hashes while their path and metadata could be held equal. Moreover, path checks,
data reading and provenance hashing occur as separate operations. No concurrent
change was performed or observed in this review, but the control flow does not
itself establish that the hashed bytes are exactly the parsed bytes under
concurrent replacement. A proposal requiring that guarantee should stage a
regular immutable input under a content identity, establish provenance from
the same captured bytes, and test its platform's path and replacement behavior.
Remote acquisition should retain provider metadata and local bytes before
analysis. These requirements strengthen the custody contract without claiming
that a real bank incident occurred.
\citep{exec23_statspai_runtime,exec23_statspai_runtime_reference}

\subsection{Re-execution and full numerical agreement are different claims}
\label{exec23:statspai-runtime-replay}
The replay helper improves on a bare call string: a standalone bundle records
a local file hash and transformations and can generate a script that checks
the input digest. Nevertheless, the generated script checks only a point
estimate when both an expected value and a result attribute exist, plus any
expected coefficient point estimates. The helper collects selected standard
error and observation-count fields into its expected-value record, but the
script does not compare them. It also lacks an assertion that at least one
number was checked, so its success text permits a zero-check case. Static
inspection of the template and independent AST assertions reproduce these
observations without generating or executing an upstream script.
This is a limited comparison contract, not proof of a numerical failure in
an estimator.
\citep{exec23_statspai_runtime,exec23_statspai_runtime_reference}

For the causal service, a replay outcome should say separately whether the
input bytes matched, whether transformations and environment matched, whether
the calculation finished, and which scientific quantities were compared.
The expected comparison set must include every claimed quantity: effect or
identified set, uncertainty representation, sample membership and weighting,
full covariance where aggregation uses it, diagnostic and refusal status,
and any random split or optimization setting relevant to the result. Missing
expected fields must fail that comparison rather than reduce its count.
Tolerances should be predeclared by statistic, with a difference attributed to
numerical error, stochastic variation or a changed target. Version text in a
script comment is not an environment lock, and agreement of a coefficient
alone cannot certify the stated uncertainty.

These findings refine the existing conditional-adapter route rather than
justify building another full estimator library. The next source review should
inspect the transformations and native fixtures actually selected for one
banking task. The later permitted prototype should then test process lifecycle,
input replacement, cache oversize handling and full replay comparisons before
scaling to more estimators. Unresolved native behavior remains a bounded
qualification obligation; commercial-service source coverage, unread frontier
papers and manuscript-depth review can proceed independently.


% SOURCE docs/release/2026-10-29/source/docs/writing-repair/2026-10-29/src/cf3018dae82d9f55.tex CHAR 0
\section{Controlled model writes preserve a conditional proof boundary}
\label{exec23:controlled-writes-scope}
A credit agent may correctly execute a rule after incorrectly reading its
premise. For example, an admitted statement that an account is in arrears
may trigger a valid collection rule even when the underlying record does
not support that statement. CPUNeSy addresses the admission boundary between
language-model proposals and a symbolic executor. Its inspected September
2026 arXiv original separates source checking, restricted predicate writes,
derivation verification and selective serving. The paper's running header
claims a 2027 conference publication; that self-description is not independently
verified acceptance as of the literature cutoff. The evidence here is the
versioned original and its explicit assumptions and reported observations.
\citep{exec23_controlled_writes_paper}

\subsection{What can be certified, and what remains a premise}
\label{exec23:controlled-writes-model}
Let $S_t$ be admitted source units, $\Gamma_t=A_t\cup R_t$ the active
facts and rules, and $\operatorname{Cn}(\Gamma_t)$ their closure in a
declared decidable Horn fragment. The source mode of the proposed gate
checks type and source support before admitting a unit. Its derived mode
requires a complete acyclic witness rooted at the claimed conclusion,
checks every edge and refuses a witness exceeding a declared depth budget.
Certificates retain dependencies and checker versions; a source erasure
requires eviction or revalidation of affected cached consequences. The
paper's state-soundness result assumes a sound source checker and a sound
fragment verifier. Its induction establishes the stated invariant under
those assumptions. It explicitly says that source soundness has not been
established for the experimental natural-language checkers.
\citep{exec23_controlled_writes_paper}

A narrower invariant is available without statistical independence. If the
allowed source vocabulary $\Omega$ excludes every rule-head predicate,
a direct write of an intermediate or target predicate is rejected. It does
not follow that every admitted leaf is true. In an independently authored
toy rule, $\texttt{arrears}\Rightarrow\texttt{deny}$, admit only
\texttt{arrears} as a writable source predicate. Directly writing
\texttt{deny} is then impossible, but a false admitted \texttt{arrears}
premise still derives it. The proposed banking contract should therefore
record both derivation status and source-support status. Neither status
establishes an intervention estimand, an adjustment set or the truth of a
borrower causal diagram.
\citep{exec23_controlled_writes_paper,exec23_controlled_writes_reference}

The voting argument likewise retains its sampling assumptions. It separates
a shared-error query mass $\beta$ from a regime with independent slips.
Conditional independence across traces is required for its concentration
bounds, including within the shared-error regime. Two agreeing calls do
not estimate $\beta$ or establish those assumptions. A stipulated example
with $\beta=1/5$, a common wrong answer on that mass and independent slip
probability $1/10$ elsewhere gives two-vote wrong, correct and abstained
probabilities $26/125$, $81/125$ and $18/125$. Answered risk is $26/107$.
Three-vote majority error is $139/625$; arbitrarily many independent votes
within each regime still leave error floor $1/5$. More calls can reduce
random slips while preserving a shared semantic mistake.
\citep{exec23_controlled_writes_paper,exec23_controlled_writes_reference}

\subsection{Execution, checking and fallback are different comparisons}
\label{exec23:controlled-writes-results}
The reported L1-HARD ablation gives 98.8\% full-pool accuracy for a one-vote
uncertified solver and 97.1\% for the restricted two-vote policy. The paper
reports no statistically significant paired difference and correctly avoids
claiming equivalence. Much of the improvement over neural controls is
attributed to deterministic execution, while certification changes the
error--abstention tradeoff. These observations do not isolate checking when
the number of grounding calls changes. A separate fixed-trace comparison
on ContractNLI supplies the more relevant check-versus-agreement evidence.
\citep{exec23_controlled_writes_paper}

For the reported DeepSeek pool of 2,091 items, agreement yields
1,533 correct, 252 wrong and 306 abstained results. Source rechecking yields
1,482 correct, 189 wrong and 420 abstained. It removes 63 wrong answers but
also withholds 51 correct answers, adding 114 abstentions. The corresponding
answered risks are $12/85$ and $63/557$, while full-pool accuracy falls by
$51/2091$. These arithmetic checks reproduce printed counts; they do not
rerun raw records or demonstrate a population guarantee. The Seed pool has
a different denominator, so the two columns are not a paired model contest.
\citep{exec23_controlled_writes_paper,exec23_controlled_writes_reference}

Fallback is a further policy. The paper's MedCalc native channel has high
answered accuracy but substantial abstention; routing withheld requests to
an uncertified same-model fallback increases full-pool accuracy. The fallback
outputs must retain their uncertified status. A CLI/MCP result should expose
which channel served each answer, the full attempted denominator, source
checker identity and uncertainty about semantic support. A single success
flag would hide the tradeoff that motivates the design.
\citep{exec23_controlled_writes_paper}

There is also a material comparator qualification inside the original.
The main public-benchmark table includes Logic-LM, DetermLR, LINC, SymbCoT,
LINA, HBLR and MenTaL rows. The appendix repeats those rows as historical
adapter summaries whose per-row split identifiers, model profiles and exact
denominators are unspecified, explicitly excluding them from primary
comparative claims. Thus the main table and its broad ranking prose cannot
support a matched superiority claim over those systems. Very low historical
scores also lack a breakdown of parsing, solver and executable-answer errors.
The comparison here retains the reported rows with this limit and draws no
ranking from their apparently common columns.
\citep{exec23_controlled_writes_paper}

\subsection{A prospective banking protocol and its cost boundary}
\label{exec23:controlled-writes-decision}
The internal hypothesis protocol has no independently verified preregistration
timestamp. Complete per-query certificate archival was not instrumented in
every adapter; model caches lack consistent sampling dates and immutable
revisions. The paper describes an offline audit package, but the inspected
PDF does not establish custody of that package in this project. Its coverage
bound is not empirically instantiated from step-level channel rates, and
independence remains unverified. Candidate-pool restriction and token savings
must also be separated from physical storage compression: reported bounded
redundancy is 37.1\% in the mathematical corpus and 0.60\% in the legal
corpus, below the internal 50\% target. Physical index-byte savings and the
scheduler-specific crossover remain unmeasured.
\citep{exec23_controlled_writes_paper}

A bounded prototype can nevertheless test a concrete question. Freeze a
public or generated banking rule corpus with explicit leaf facts, source
spans, admissible predicates and known interventions before held-out queries.
Compare an unrestricted translator plus solver, restricted writes, fixed-trace
source rechecking and a separately labeled fallback. Inject both prohibited
target writes and false but well-typed leaves. Test dependency deletion,
changed-source cache invalidation, missing evidence, contradictory records,
abstention and full failure reporting. Independent labels must judge source
faithfulness separately from formal derivability. This is a prospective
proposal, not an implemented source checker or live-model trial.
\citep{exec23_controlled_writes_paper}

For a transparent economic decision, let $c_W$ price a wrong served answer,
$c_A$ an abstention and $C$ the additional total checking cost. Using the
printed DeepSeek changes, loss savings are $63c_W-114c_A-C$. Stipulating
$c_W=10$ dollars and $C=100$ dollars gives USD 302 at $c_A=2$ and
USD $-154$ at $c_A=6$. The change of sign illustrates why accepted-answer
precision alone is not an investment criterion. The costs are invented
inputs for the displayed example, explicitly editable in the independent
reference; no bank value or causal performance gain is estimated.
\citep{exec23_controlled_writes_reference}

\section{What the complete evaluation must preserve}

The evaluation argument now reaches from the decision to its operational evidence. A benchmark needs the right causal target and a trustworthy reference before accuracy has a useful meaning. Exposure, selective adjudication and repeated calls alter what a measured score can support. Stability, agreement, derivability and successful execution each test a narrower property; their counterexamples explain why none can replace scientific correctness. The toolkit studies add an equally concrete requirement: the uncertainty, failures and input identity must survive the interface through which the result is evaluated.

These restrictions still permit a useful first experiment. Admit bounded public or generated problems with independent references, compare declared information and resource budgets, retain all assigned outcomes, and evaluate decision loss alongside availability and conditional accuracy. Choose the sample size from the particular claims and plausible effect margins, not from a convenient collection. A live comparison then addresses the remaining translation and review costs; an institutional study addresses the bank-specific assignment, transfer and monetary questions. Each stage can justify the next without borrowing a guarantee from a different target, score or software layer.

